% !TEX root = ../../main.tex
\setchapterimage{figure/chapter1封面图.jpg}
\setchapterintro{}
\chapter{整数与多项式理论}
\section{整数理论}
根据我们在分析学笔记里所做的工作，$\Z$与$\Q$已经被我们从无到有的构建出来.我们曾经指出，$\Z$是一个整环,而 $\mathbb{Q}$ 构成一个数域，下面，我们对代数结构的概念进行重申.
\begin{definition}\phantomsection\label{p1:concept:c1-l7}
	代数结构，是一个配备了一种或多种运算的多元组，下面的初级的代数结构中，用＂$\oplus$＂来表示其合成法则
	\begin{itemize}
		\item 半群$S _\oplus, \forall a, b,c \in S _\oplus$ ，满足
		\begin{enumerate}
			\item$a \oplus b \in S _\oplus$
			\item$(a \oplus b) \oplus c=a \oplus(b \oplus c)$
		\end{enumerate}
		\item 幺半群$M _\oplus \quad \forall a, b,c \in M _\oplus$ 满足
		\begin{enumerate}
			\item $a \oplus b \in M _\oplus$
			\item$(a \oplus b) \oplus c=a \oplus(b \oplus c)$
			\item$\exists i \in M _\oplus$ s．t．$i \oplus a=a \oplus i=a$
		\end{enumerate}
		\item 群 $G_ \oplus \quad \forall a, b, c \in a \oplus \quad$ 满足
		\begin{enumerate}
			\item $a \oplus b \in G _\oplus$
			\item$(a \oplus b) \oplus c=a \oplus(b \oplus c)$
			\item$\exists i \in G _\oplus$ s．t．$i \oplus a=a \oplus i=a$
			\item $\exists b \in G _\oplus$ s．t．$a \oplus b=b \oplus a=i$
		\end{enumerate}
		\item Abel群 $G_ \oplus \quad \forall a, b, c \in a \oplus \quad$ 满足
		\begin{enumerate}
			\item $a \oplus b \in G _\oplus$
			\item$(a \oplus b) \oplus c=a \oplus(b \oplus c)$
			\item$\exists i \in G _\oplus$ s．t．$i \oplus a=a \oplus i=a$
			\item $\exists b \in G _\oplus$ s．t．$a \oplus b=b \oplus a=i$
			\item $a\oplus b=b\oplus a$
		\end{enumerate}
		下面的代数结构中有两种运算，一种叫加法，记为＂$+$＂一种叫乘法，记为＂$\times$＂
		
		
		
		\item 环$R$
		
		$R_{+}$是Abel群,$R_\times$ 是半群或幺半群.
		\item 交换环$R$
		
		$R_{+}$是Abel群,$R_\times
		$ 是交换半群或交换幺半群.
		\item 整环
		
		$R_{+}$是Abel群,$R_\times$ 是交换幺半群，且无零因子( 零因子是指$a, b \neq 0$ s．t．$a b=0$)
		\item 域$\F$
		
		$R_{+}$是Abel群,$R_\times$ 是Abel群，且无零因子( 零因子是指$a, b \neq 0$ s．t．$a b=0$)
	\end{itemize}
	
	
\end{definition}

下面，我们讨论整数理论.
\begin{definition}\phantomsection\label{p1:concept:c1-l59}
	$a, b \in \mathbb{Z}$ ，若 $\exists q \in \mathbb{Z}\text{且}q\ne0$ s．t．$a=b q$  ,则说 $b$ 整除 $a$ 或 $b$ 是 $a$ 的因子，用 $b \mid a$ 表示，我们也会使用记号 $b \nmid a$ 表示 $b$ 不整除$a$ .
	
	对 $a \neq 1$, 我们称 $1,-1, a,-a$ 为 $a$ 的平凡因子.
	
	特别的，对 $\forall a \in \mathbb{Z}$ ，都有 $a\mid0$.
	
\end{definition}

\begin{proposition}整除的性质
	
	
	\begin{enumerate}
		\item$a|b \Longleftrightarrow-a| b \Longleftrightarrow a|-b \Longleftrightarrow| a|||b|$ ．
		\item$a \mid b$ 且 $b|c \Rightarrow a| c$ ．
		\item$a \mid b$ 且 $a \mid c \Longleftrightarrow$ 对任意的 $x, y \in \Z$ ，有 $a \mid b x+c y$ ．
		
		
		一般地，$a\left|b_1, \cdots, a\right| b_k$ 同时成立 $\Longleftrightarrow$ 对任意的 $x_1, \cdots, x_k \in \Z$ ，有 $a \mid b_1 x_1+\cdots+b_k x_k$.
		\item 设 $m \neq 0$ ，那么 $a|b \Longleftrightarrow m a| m b$ ．
		\item$a \mid b$ 且 $b \mid a \Longrightarrow b= \pm a$ 
		\item 设 $b \neq 0$ ，那么 $a|b \Longrightarrow| a|\leqslant|b|$ ．
		\item 设 $a \neq 0, b=q a+c$ ，那么 $a$ 整除 $b$ 的充分必要条件是 $a$ 整除 $c$ .
	\end{enumerate}
\end{proposition}
\begin{proof}
	\begin{enumerate}
		\item$b=a q \Leftrightarrow b=(-a)(-q) \Leftrightarrow-b=a \cdot(-q) \Leftrightarrow|b|=|a||q|$
		\item$b=a q_1, c=b q_2 \Leftrightarrow c=\left(q_1 q_2\right) a$
		\item$b=k_1 , c=k_2 a \Rightarrow b x+c y=\left(k_1 x+k_2 y\right) a \Rightarrow a \mid b x+c y$,对推论同理
		\item$b=a q \Leftrightarrow m b=m a  q\Leftrightarrow m a| m b$
		\item $a=q b, b=q^{\prime} a \Rightarrow 1=q q^{\prime} \Rightarrow q=q^{\prime}= \pm 1 \Rightarrow a= \pm b$
		\item $a|b \Leftrightarrow| a|||b| \quad| b|=|q||a| \quad$ 因 $q \in \mathbb{Z} \quad$ 且 $b \neq 0, q \neq 0 \Rightarrow|b| \geqslant|a|$
		\item $\Rightarrow: a \mid b \Rightarrow b=k a \Rightarrow k a=q a+c \Rightarrow c=(k-q) a\Rightarrow a \mid c$,
		
		$\Leftarrow :a|c \Rightarrow c=r a \Rightarrow b=(q+r) a \Rightarrow a| b$
		
	\end{enumerate}
\end{proof}
\begin{definition}\phantomsection\label{p1:concept:c1-l98}
	称 $p$ 为素数，若 $p$ 只有平凡因子
\end{definition}
\begin{definition}\phantomsection\label{p1:concept:c1-l101}
	若 $c \mid a$ 又 $c \mid b$ ，则称 $c$ 为 $a,b$ 的公因子，记$a,b$所有公因子构成的集合为$\mathscr{D}(a,b)$,若 $d$ 为 $a, b$ 所有公因子里最大的那个，也就是说$d=\max\mathscr{D}(a,b)$就称 $d$ 为 $a, b$ 的最大公因子，记为 $d=(a, b)$.
	
	若 $(a, b)=1$ ，称 $a, b$ 互素.
	
	对 $a_1, \cdots, a_n$ ，它们的公因数写成的集合记作$\mathscr{D}(a_1,\cdots,a_n)$，最大公因数$d=\max\mathscr{D}(a_1, \cdots, a_n)$,记作 $\left(a_1, \cdots, a_n\right) $.
	
	
	特别地 $(a, 0)=|a|$ 而 $(0,0)$ 无定义.
\end{definition}
\begin{theorem}
	$ a_1, \cdots, a_n \in \mathbb{Z}$
	
	$$
	\left(a_1, a_2, \cdots, a_n\right)=\left(\left(a_1, a_2\right), \cdots, a_n\right)
	$$
	
	
	
\end{theorem}
\begin{proof}
	设 $d=\left(a_1, a_2, \cdots, a_n\right) , h=\left(a_1, a_2\right) , g=\left(h, a_3, \cdots, a_n\right)$
	
	$ d=\left(a_1, a_2, \cdots, a_n\right) \Rightarrow d\mid a_1, d\mid  a_2 \Rightarrow d \mid h$，同时 $ d \mid a_3, \cdots, a_n$, 故 $d \mid g $
	
	$g \mid h$ 而 $h \mid a_1$ 和 $a_2$ ，故 $g \mid a_1, a_2$ ,同时 $g \mid a_3, \cdots, a_n$,故 $g \mid d$
	
	园此 $g \mid d$ 且 $d \mid g$,由最大公因数的定义 $g=d$.
\end{proof}

\begin{proposition}
	\begin{enumerate}
		\item$	\left(a_1, a_2\right)=\left(a_2, a_1\right)=\left(-a_1, a_2\right)=\left(\left|a_1\right|,\left|a_2\right|\right) $
		
		一般地，有
		
		$$
		\begin{aligned}
			& \left(a_1, a_2, \cdots, a_i, \cdots, a_k\right)=\left(a_i, a_2, \cdots, a_1, \cdots, a_k\right) \\
			& \quad=\left(-a_1, a_2, \cdots, a_k\right)=\left(\left|a_1\right|, \cdots,\left|a_i\right|, \cdots,\left|a_k\right|\right)
		\end{aligned}
		$$
		\item 若 $a_1 \mid a_j, j=2, \cdots, k$ ，则
		
		$$
		\left(a_1, a_2\right)=\left(a_1, a_2, \cdots, a_k\right)=\left|a_1\right| .
		$$
		\item 对任意的整数 $x,\left(a_1, a_2\right)=\left(a_1, a_2, a_1 x\right)$ ，
		
		$$
		\left(a_1, \cdots, a_k\right)=\left(a_1, \cdots, a_k, a_1 x\right)
		$$
		
		\item 对任意整数 $x,\left(a_1, a_2\right)=\left(a_1, a_2+a_1 x\right)$ ，
		
		$$
		\left(a_1, a_2, a_3, \cdots, a_k\right)=\left(a_1, a_2+a_1 x, a_3, \cdots, a_k\right)
		$$
		
		\item 若 $p$ 是素数，则
		
		$$
		\left(p, a_1\right)= \begin{cases}p, & p \mid a_1, \\ 1, & p \nmid a_1 .\end{cases}
		$$
		
		
		一般地，有
		
		$$
		\left(p, a_1, \cdots, a_k\right)= \begin{cases}p, & p \mid a_j, j=1,2, \cdots, k, \\ 1, & \text { 其他. }\end{cases}
		$$
		
		
	\end{enumerate}
	
	
\end{proposition}
\begin{proof}
	根据公因数的定义及整除性质推出
	
	$$
	\begin{aligned}
		& \mathscr{D}\left(a_1, a_2\right)=\mathscr{D}\left(a_2, a_1\right)=\mathscr{D}\left(-a_1, a_2\right)=\mathscr{D}\left(\left|a_1\right|,\left|a_2\right|\right), \\
		& \mathscr{D}\left(a_1, a_2\right)=\mathscr{D}\left(a_1, a_2, a_1 x\right), \quad x \in \Z, \\
		& \mathscr{D}\left(a_1, a_2\right)=\mathscr{D}\left(a_1, a_2+a_1 x\right), \quad x \in \Z .
	\end{aligned}
	$$
	
	
	由此及最大公因数的定义就分别证明了1,3,4当 $k=2$ 时成立，$k>2$ 的情形同样证明.2由 1.1.3 的6推出．5由素数的定义及2推出．证毕．
\end{proof}
\begin{theorem}
	$a, b \in \mathbb{Z}\textbf{且} b \neq 0$ ，$\exists$ 唯一 $q, r \in \mathbb{Z}, 0 \leq r<|b|$ 使得 $a=b q+r$
	成立.我们叫这种算法带余除法.
	
\end{theorem}
\phantomsection\label{p1:concept:c1-l197}一个满足带余除法的环，叫做欧几里德环.
\begin{proof}
	唯一性：设 $a=b q_1+r_1=b q_2+r_2 \quad r_{1,2}=0$ 或 $0<r_{1,2}<|b|$
	
	那么 $b\left(q_1-q_2\right)=r_2-r_1 \Rightarrow b \mid r_2-r_1$ 或 $r_2-r_1=0$
	若 $b \mid r_2-r_1$ ,则由 1.1.3 (6)有 $|b| \leqslant\left|r_2-r_1\right|$而$\left|r_2-r_1\right|<|b|$ 矛盾
	
	故$r_2-r_1=0 \quad r_2=r_1$ 且 $q_1=q_2 \quad$ 
	
	唯一性得证.
	
	存在性：$S=\{a-t b: t \in \mathbb{Z}\}$
	
	令 $r$ 是 $S$ 中最小的非零元素
	
	$\exists q \in \mathbb{Z}$ s．t $r=a-q b$ 也就是 $a=q b+r$
	
	
	断言 $r=0$ 或 $0<r<|b|$
	
	
	$r=0$ 时，这是平凡的.
	
	$r>0$ 时，若 $0<r<|b|$ 不成立，设 $r \geqslant|b|$ 不妨设 $b>0$
	
	$$
	r>b \Rightarrow r-b>0 \Rightarrow a-(q+1) b \in S
	$$
	
	
	而 $r>r-b$ 与 $r$ 的取法矛盾，
	
	若 $b<0 \Rightarrow r>-b \Rightarrow r+b>0 \Rightarrow a-(q-1) b \in S$
	
	$r>r+b \quad$ 同样与 $r$ 的取法矛盾
	
	唯一性得证.
\end{proof}
\begin{problem}
	二进制与十进制有如下转换规律
	若 $\left(a_n a_{n-1} \cdots a_2 a_1 a_0\right)_2=[n]_{10} \quad a_n \cdots, a_0 \in\left\{1,0\right\}$
	
	$$
	n=a_0+a_1 \cdot 2+a_2 \cdot 2^2+\cdots+a_n 2^n \quad a_n=1
	$$
	
	
	给出一个从 10 进制数转换 2 进制的算法.算$(2001)_{2}$的二进制数.
	
	
\end{problem}
\begin{solution}
	对 $n$ 不断除 2 ，取其余数
	
	$\begin{aligned}
		\textbf{商数}\quad&\textbf{除数}&\textbf{余数}\\
		a_0+a_1 \cdot 2+a_2 \cdot 2^2 \cdots+a_n 2^n\quad&2&a_0\\
		a_1+a_2 \cdot 2+\cdots+a_n \cdot 2^{n-1} \quad &2&a_1\\
		a_2+a_3^2+\cdots+a_n \cdot 2^{n-2} \quad&2&a_2\\
		\vdots\\
		\quad&2&a_n
	\end{aligned}$
	
	这样$(n)_{10}=\left(a_n a_{n-1} \cdots a_2 a_1 a_0\right)_2$
	
	对2001
	
	$\begin{aligned}
		\textbf{商数}\quad&\textbf{除数}&\textbf{余数}\\
		2001\quad&2&1\\
		1000\quad&2&0\\
		500 \quad &2&0\\
		250\quad&2&0\\
		125\quad&2&1\\
		62\quad&2&0\\
		31\quad&2&1\\
		15\quad&2&1\\
		7\quad&2&1\\
		3\quad&2&1\\
		1\quad&2&1\\
	\end{aligned}$
	
	$2001=(11111010001)_2$
\end{solution}
\begin{lemma}
	$ a, b \in \mathbb{Z} \quad b \neq 0 \text { 必有 } a=b q+r \Rightarrow(a, b)=(b, r)$
\end{lemma}
\begin{proof}
	只需证$\mathscr{D}(a,b)=\mathscr{D}(b,r)$
	
	1.$\mathscr{D}(a,b)\subset\mathscr{D}(b,r)$
	
	$\forall d \in \mathbb{Z}$ 且 $d|a, d| b$ ，有 $a=k_1 d , b = k_2 d$ 其中 $k_1, k_2 \in \mathbb{Z}$
	
	$ r=a-b q=\left(k_1-q k_2\right) d \Rightarrow d \mid r$ 又 $d \mid b$
	
	故$\mathscr{D}(a,b)\subset\mathscr{D}(b,r)$
	
	2.$\mathscr{D}(b,r)\subset\mathscr{D}(a,b)$
	
	$\forall d \in \Z$ 且 $d\mid b, d\mid r$ 有 $b=k_1 d , r=k_2 d$ 其中 $k_1, k_2 \in\Z$
	
	$ a=b q+r=\left(k_1 q+k_2\right) d \Rightarrow d \mid a $且$ d \mid b$
	
	故$\mathscr{D}(b,r)\subset\mathscr{D}(a,b)$
	
	综上$\mathscr{D}(a,b)=\mathscr{D}(b,r)$.
\end{proof}

这样我们就得到求$(a, b)$  的一种算法，欧几里德算法.


\begin{algo}欧几里得算法
	
	$\begin{aligned}
		& a=b q_1+r \\
		& b=r q_1+r_1 \\
		& r_1=r_2 q_2+r_3 \\
		& \vdots \\
		& r_{k-2}=r_{k-1} q_k+r_k\\
		&r_k=0\textbf{且}r_{k-1}\mid r_{k_2}
	\end{aligned}$
	
	由引理1.1.9
	$$(a, b)=(b, r)=\left(r, r_1\right)=\cdots=\left(r_{k-2}, r_{k-1}\right)=\left(r_{k-1}, r_k\right)=r_{k-1}$$
\end{algo}
\begin{definition}
	$a, b \in \mathbb{Z} \quad \exists u, v \in \mathbb{Z}$ s．t
	
	$$
	(a, b)=a u+b v
	$$
	
	
	
\end{definition}
\phantomsection\label{p1:concept:c1-l333}这称为辗转除法.
\begin{proof}
	法一：由 $1.1.9 \quad d=(a, b)=r_{k-1}$
	
	$$
	\begin{aligned}
		& r_{k-1}=r_{k-3}-q_{k-1} r_{k-2} \\
		& r_{k-2}=r_{k-4}-q_{k-2} r_{k-3}\\
		&\vdots\\
		& r_1=b-q_1 r \\
		& r=a-b q
	\end{aligned}
	$$
	
	
	那么
	
	$$
	\begin{aligned}
		(a, b)=r_{k-1} & =r_{k-3}-q_{k-1}\left(r_{k-4}-q_{k-2} r_{k-3}\right) \\
		& =\left(1+q_{k-1} q_{k-2}\right) r_{k-3}-q_{k-1} r_{k-4}
	\end{aligned}
	$$
	
	
	迭代下去，我们得到 $a, b$ 和 $u, v$.
	
	法二:构造 $\mathcal{X}=\{z: z=x a+y b,  x,  y \in \mathbb{Z}\}$ ，令 $d$ 为$\mathcal{X}$中最小的非负整数.
	
	下证 $d=(a, b)$
	
	$$\exists u,v\in\Z,
	d=u a+v b
	$$
	
	
	现证 $d \mid a$ ,设 $a=q d+r $ ,.$ r=0$ 或 $0<r<d $
	
	$r=0$时， $d\mid a$显然.
	
	
	$0<r<d \textbf{时,} r=a-q d=a-(u a+u b) q \in \mathcal{X}$
	
	$r<d$ 与 $d$ 取法矛盾 故 $d \mid a \quad$ ,同理 $d \mid b$.
	
	对 $d_1\mid a, d_1\mid b$ 有 $d_1\mid a u+b v \Rightarrow d_1\mid d$
	
	综上 $(a, b)=d$
\end{proof}

这样，我们给出 $a, b$ 最大公因子的第二定义.

\begin{definition}
	$a, b$ 的最大公因子是集合$\left\{z:z= x a+y b, x, y \in \Z\right\} $ 中最小的非负整数.
	
\end{definition}
\begin{problem}
	$a_1 ,\cdots, a_n \in \mathbb{Z} \quad d \in \mathbb{Z}$
	那么证：$\left(d a_1, \cdots, d a_n\right)=|d|\left(a_1, \cdots, a_n\right)$
\end{problem}
\begin{solution}
	设 $\left(d a_1, \cdots, d a_n\right)=m,\left(a_1, \cdots, a_n\right)=l$
	
	那么 $\exists u_1, \cdots, u_n \in \mathbb{Z}$, s．t． $ l=\sum_{i=1}^{n} u_ia_i$
	
	$|d| l=u_1 |d|a_1+\cdots+u_n |d| a_n$，
	由最大公因数的第二定义，$l=\Sigma u_i a_i$ 是 $\forall k_1, \cdots k_n$  $\sum k_i a_i$ 中最小的非负整数.
	
	$|d| l$就是 $\sum k_i|d|a_i$中最小的非负整数，也就是$m$.
	
	综上，	$|d|\left(a_1, \cdots, a_n\right)=\left(d a_1, \cdots, d a_n\right)$
	
\end{solution}
\begin{remark}
	此处的 $u, v$ 不是唯一的.
	
	对 $u a+v b=d $, 有 $(u-t b) a+(v-t a) b=d \quad t \in \mathbb{Z}$.
\end{remark} 
\begin{corollary}
	$ a_1, \cdots, a_s \in \mathbb{Z}$ 不全为 $0, \exists u_1, \cdots, u_s \in \mathbb{Z}$s.t.
	
	$$
	\sum_{i=1}^s u_i a_i=\left(a_1, \cdots, a_s\right)
	$$
\end{corollary}
\begin{proof}
	$s=2$ 时 成立
	
	设 $s=3, \cdots, n-1$ 成立
	
	先证
	$$
	\left(\left(a_1, \cdots, a_{n-1}\right), a_n\right)=\left(a_1, \cdots, a_n\right)
	$$
	
	设 $\left(a_1, \cdots, a_n\right)=d$ ,$\left(a_1, \cdots, a_{n-1}\right)=h$ ,$ \left(h, a_n\right)=g$
	
	
	$d\left|a_1, \cdots, a_n \Rightarrow d\right| a_n$ 且 $d\left|a_1, \cdots, a_{n-1} \Rightarrow d\right| a_n$ 且 $d|h \Rightarrow d| g$
	
	
	$g \mid h$ 且 $h\left|a_1, \cdots, a_{n-1} \Rightarrow g\right| a_1, \cdots, a_{n-1}$ 且 $g\left|a_n \Rightarrow g\right| d$
	
	$$
	d=g \Rightarrow\left(\left(a_1, \cdots, a_{n-1}\right), a_n\right)=\left(a_1, \cdots, a_n\right)
	$$
	
	对 $\left(a_1, \cdots, a_{n-1}\right)=d_{n-1} \quad \exists u_1, \cdots, u_{n-1}$ s．t．$d_{n-1}=\sum_{i=1}^{n-1} u_i a_i$
	
	对 $\left(d _{n-1}, a_n\right)=d $,$ \exists v_1, v_2$ s．t．
	
	$$\begin{aligned}
		d & =v_1 d_{n-1}+v_2 a_n \\
		& =\sum_{i=1}^n t_i a_n
	\end{aligned}$$
	
	$t_i=v_1 u_i \quad 1 \leqslant i \leqslant n-1 , t_n=v_2$
	
	
\end{proof}
\begin{corollary}
	$p$ 是素数，$p \mid a b,(a, b)=1$ 那么 $p \mid a$ 或 $p \mid b$ .
\end{corollary}
\begin{proof}
	若 $p \nmid a\Rightarrow(p, a)=1, \exists u, v \in \mathbb{Z}$ s．t．$u p+v a=1$
	
	且 $\exists q$ s．t．$p q=a b \Rightarrow u p+v a=1 \Rightarrow u p b+v a b=b \Rightarrow p(u b+v q)=b$
	$\Rightarrow p \mid b$
	
	$p\nmid b \Rightarrow p \mid a$ 同理
	
\end{proof}

\begin{theorem}算术基本定理
	
	任何一个非零整数 $a$ ，都可以写成唯一的形式
	
	$$
	a=(-1)^{\varepsilon(a)} p_1 \cdots p_n
	$$
	
	$\varepsilon(a)$ 是 $a$ 的符号指标 $\varepsilon (a)=\left[\begin{array}{ll}0 & a>0 \\ 1 & a<0\end{array}\right.$ 
	
	$p_1, \cdots, p_n$ 是素数，可以相同.这里 $a\ne \pm1$
	
\end{theorem}
\begin{proof}
	存在性:若存在一个整数不能够分解为如上形式,设不能分解为上述形式的最小整数为 $n$ .
	
	若 $n$ 为素数 得证. 若 $n$ 不是素数，那么存在 $1<n_1, n_2<n $ s．t．$n=n_1n_2$而 $n_1, n_2$ 可分解得 $n$ 和可分解，矛盾.
	
	唯一性 ：设 $a=p_1 p_2 \cdots p_1=q_1 q_2 \cdots q_s$
	由 1.1.16 ，$\exists$唯一 $i$ 使 $p_1 \mid q_i$ ，适当调序使 $i=1 $这样$ p_1 \mid q_1$.
	
	因 $p_1, q_1$ 是素数，有
	$$p_1=\left(p_1, q_1\right)=q_1$$
	
	由消去律 $p_2 p_3 \cdots p_r=q_2 q_3 \cdots q_s$
	
	同样的方式 $p_2=q_2, p_3=q_3 \cdots$ 最后 $\quad r=s$
\end{proof}
这样我们给出$a$标准分解式
$$
a=(-1)^{\varepsilon(a)} \prod_{i=1}^{n}p_1^{r_i}
$$
\phantomsection\label{p1:concept:c1-l498}$p_1, \cdots, p_n$ 不相同．记 $n_i=\text{ord}_{p_i}(a)$ 称为 $a$ 中 $p_i$ 的阶数.
\begin{theorem}
	素数是无穷多的.
	
	
	
	
\end{theorem}
\begin{proof}
	如若不然，设 $p_1 \because p_n$ 是所有的素数.
	
	$\text { 令 } N=p_1 \cdots p_n+1>\max(p_n, \cdots, p_1 )\quad \text { 对 } N \quad \exists p_i \text { s.t. } N=k p_i\Rightarrow1=p_1\cdots p_i\cdots p_n-kp_i\Rightarrow p_i\mid1
	$
	矛盾
	
	故素数有无穷多.
	
\end{proof}
\section{数域，扩域与子域}
我们对域的定义进行重申.
\begin{definition}\phantomsection\label{p1:concept:c1-l518}
	对$(\F,+,\cdot)$
	\begin{enumerate}
		\item 对四则运算封闭
		
		\item 满足交换律，结合律，分配律
	\end{enumerate}
\end{definition}
\begin{definition}\phantomsection\label{p1:concept:c1-l526}
	设 $\F$ 和 $\E$ 是两个域，并且 $\F$ 是 $\E$ 的子集（ $F \subseteq E$ ）.如果 $\E$ 中的加法和乘法运算，限制在 $\F$ 上时，与 $\F$ 原来的加法和乘法运算一致，那么我们称：
	\begin{itemize}
		\item  $E$ 是 $F$ 的一个扩域.
		\item $F$ 是 $E$ 的子域.
		\item 这个扩域关系记作 $E / F$.注意，这里的斜线不是除号.
	\end{itemize}
	
	
	
\end{definition}


\begin{proposition}
	$\Q$是最小的数域.
\end{proposition}
\begin{proof}
	$\Q$ 是最小的数域$\Leftrightarrow$ 所有的数域$\F$包含$\Q$作为子域
	
	设 $a \in \mathbb{F} z \quad a / a=1$ ，由 1 统对四则运算的封闭，生成了 $\Q\mid\left\{0\right\}$ ，又$ a-a=0$, 这样就生成了 $\Q$
	那么任意的一个数域$\F$都包含$\Q$作为子域.
	
\end{proof}

\begin{example}
	\[\mathbb{Q}(\sqrt{2})=\{z: z=a+b \sqrt{2}, a, b \in \mathbb{Q}\} \textbf{是包含} \sqrt{2}\textbf{的最小数域.} 
	\]
	
	\[\mathbb{Q}(\pi)=\left\{z: z=\frac{a_n \pi^n+\cdots+a_0}{b_m \pi^m+\cdots+b_0}, a_n \cdots, a_0, b_m, \cdots, b_0 \in \mathbb{Q}, b_m \cdots, b_0\textbf{不全为0}\right\}\]
	是包含$\pi$的最小数域.		
\end{example}
\section{一元多项式}
\begin{definition}\phantomsection\label{p1:concept:c1-l558}
	设函数  $f(x) = a_n x^n + a_{n-1} x^{n-1} + \cdots + a_1 x + a_0 $，其中 $ a_i \in \mathbb{F}$, $i=0,1,\cdots,n$ .
	称  $f $ 为一个定义在数域$\F$上的多项式，其中  $x $ 是不定元.
	在数域  $\mathbb{F}$  上所有多项式构成的集合称为  $\mathbb{F} $ 上的多项式环，记作 $ \mathbb{F}[x]$ .
	\begin{itemize}
		\item 称  $a_i (i=0,1,\cdots,n)$  为  $f$  的系数.若  $a_n \neq 0$ ，则称 $a_n$  为首系数，$ a_n x^n $ 为首项，$n $ 称为  $f(x)$ 的次数，记作 $ \partial f = n $ 或  $\deg f = n$ .
		\item 若  $f(x) = c \neq 0 $，则  $\deg f = 0$  且反之亦然， $\deg f = 0$  则 $ f(x) = c \neq 0$ .
		\item 若 $ a_n = 1$ ，则  $f(x)$  称为首一的.
		\item 零多项式 $f(x) = 0$  的阶数规定为 $-\infty$ .
		\item $\forall f,g\in \F[x]$,若$f,g$的所有系数相等，则$f=g$
		\item $\forall f,g\in \F[x]$若$f=cg,c\in\F\wedge c\neq0$,称$ f,g$相伴，显然地，相伴关系是等价关系.
	\end{itemize}
	
\end{definition}
可以把多项式$f(x)= a_n x^n + a_{n-1} x^{n-1} + \cdots + a_1 x + a_0$理解为无穷维空间上的点$(a_0,a_1.\cdots,a_n,0,0,\cdots)$.
\begin{definition}\phantomsection\label{p1:concept:c1-l573}
	对于 \( f, g \in \mathbb{F}[x] \)，定义：
	
	\[
	f = a_n x^n + \cdots + a_0 = \sum_{i=0}^{n} a_i x^i
	\]
	
	\[
	g = b_m x^m + \cdots + b_0 = \sum_{i=0}^{m} b_i x^i
	\]
	
	其中 \( a_n, b_m \neq 0 \)，且 \( \deg f = n \), \( \deg g = m \).不妨设 \( m > n \).将 \( f \) 的系数向 \( g \) 的系数补齐，设
	
	\[
	a_m = \cdots = a_{n+1} = 0 
	\]
	
	那么定义 \( f + g \) 为：
	
	\[
	f + g := \sum_{i=0}^{m} (a_i + b_i) x^i
	\]
	类似地，定义 \( f - g \) 为：
	
	\[
	f - g := \sum_{i=0}^{m} (a_i - b_i) x^i
	\]
	
	定义\( f ,g \) 的乘法：
	
	\[\begin{aligned}
		fg &:= a_n b_m x^{n+m} + (a_{n-1} b_m + a_n b_{m-1}) x^{m+n-1} + \cdots + a_0 b_0\\
		&= \sum_{k=0}^{m+n} \left( \sum_{i+j=k} a_i b_j \right) x^k
	\end{aligned}\]
	
\end{definition}
显然有：
	\[\deg(fg) = \deg f + \deg g\]
	\[\deg(f + g) \leqslant \max \{ \deg f, \deg g \}\]
	
多项式的二元运算满足以下规律
\begin{proposition}
	对于 \( f, g, h \in \mathbb{F}[x] \)，有以下性质：
	
	\begin{itemize}
		\item[A1.] 加法结合律：\( (f + g) + h = f + (g + h) \)
		\item[A2.] 加法交换律：\( f + g = g + f \)
		\item[A3.] 加法单位元：\( \exists 0 \in \mathbb{F}[x] \)，使得 \( f + 0 = f \)
		\item[A4.] 加法逆元：\( \forall f \in \mathbb{F}[x] \)，\( \exists g \in \mathbb{F}[x] \)，使得 \( f + g = 0 \)
		\item[M1.] 乘法结合律：\( (fg)h = f(gh) \)
		\item[M2.] 乘法交换律：\( fg = gf \)
		\item[M3.] 乘法单位原：\( \exists 1 \in \mathbb{F}[x] \)，使得 \( f \cdot 1 = f \)
		\item[D1.] 分配律：\( f(g + h) = fg + fh \)
		\item[D2.] 消去律：\( fg = fh \)，则 \( g = h \) 或 \( f = 0 \)
	\end{itemize}

\end{proposition}
	
证明是容易的，由此可见，\( \mathbb{F}[x] \) 是一个环.

与整数环类似的，\( \mathbb{F}[x] \)上也有整除关系.
\begin{definition}\phantomsection\label{p1:concept:c1-l634}
	\( f, g \in \mathbb{F}[x] \)，若存在 \( h \in \mathbb{F}[x] \)，满足
	\[ f = h \cdot g \]
	则称 \( g \) 整除 \( f \) 或 \( g \) 是 \( f \) 的因子，记作 \( g \mid f \).若 \( g \) 不整除 \( f \) 则记作 \( g \nmid f \).
	显然地，阶数为$0$ 的多项式是任意 \( f \) 的因子，并且任何非零多项式是 $0$ 多项式的因子.若 \( g \mid f \)，则 \( \deg g \leq \deg f \).
	
\end{definition}
显然地，若 \( f \mid g \) 且 \( f \mid h \)，则
\[ f \mid (gu + hv) \quad \forall u, v \in \mathbb{F}[x]. \]

\phantomsection\label{p1:concept:c1-l644}若 \( d \mid f \) 且 \( d \mid g \)，则称 \( d \) 是 \( f, g \) 的公因子.若 \( f, g \) 的任意公因子都可从被 \( d \) 整除，则称 \( d \) 为 \( f, g \) 的最大公因式.记为 \( (f, g) \)，与整数环上类似，多项式环上也有欧几里得算法.若 \( (f, g) = 1 \)，则称 \( f, g \) 互素.

\( (f(x), 0) = cf \), \( c \) 是非零且使$f$成为首一多项式

与整数环类似，多项式环上也有带余除法和欧几里得算法
\begin{proposition}(带余除法)\label{带余除法}
	 \( \forall f, g \in \mathbb{F}[x] \)，且 \( g \neq 0 \)，存在唯一 \( q, r \in \mathbb{F}[x] \)，满足 \( r = 0 \) 或 \( \deg r < \deg g \)，使得
	\[ f = qg + r. \]
	
	
\end{proposition}
\begin{proof}
	先证明唯一性.假设
	\[ f = q_0 g + r_0=qg + r, \]
	其中 \( r(x) = 0 \) 或 \( \deg r < \deg g \)，\( r_0(x) = c \)或 \( \deg r_0 < \deg g \).
	
	那么
	\[ (q - q_0)g = r_0 - r. \]
	
	其中，\( r_0 - r = 0 \) 或 \( \deg(r_0 - r) < \deg g \)，故 \( g \mid r_0 - r \).
	
	这意味着
	\[ r(x) = r_0(x) \]
	
	因此，
	\[ q(x) = q_0(x)\]
	
	下面证明存在性.设 \( \deg f = m \)，\( \deg g = n \)，若$m<n$,我们取$q=0\wedge r=f$.若 \( m \geq n \)，我们对$\deg f$进行归纳.
	
	令$a_m x^m$与$b_nx^n$为$f$与$g$的首项，那么
	\[ f(x) = f_1(x) +\frac{a_m}{b_n} x^{m-n} g(x), \quad \deg f_1 < \deg f \text{ 或 } f_1(x) = 0. \]
	
	若
	\[ f_1(x) = 0\]
	取
	\[ q(x) = \frac{a_m}{b_n} x^{m-n}, \quad r(x) = 0. \]
	
	否则，
	\[ \deg f_1 < \deg f. \]
	
	由归纳假设，存在 \( q_1 \) 和 \( r \) 满足 \( r = 0 \) 或 \( \deg r < \deg g \)，使得
	\[ f_1 = q_1 g + r. \]
	
	因此，
	\[ f = qg + r, \]
	其中
	\[ q = \frac{a_m}{b_n} x^{m-n} + q_1, \]
	且
	\[ r = 0 \text{ 或 } \deg r < \deg g. \]
	
	这样，证明了 \( q(x) \) 和 \( r(x) \) 的存在性.
	
\end{proof}

使用带余除法我们可以得到与整数环类似的欧几里得算法

\begin{proposition}(多项式环上的欧几里得算法)\label{多项式环上的欧几里得算法}
对于任意 \( f, g \in \mathbb{F}[x] \)，设 \( d = (f, g) \)，则存在 \( u, v \in \mathbb{F}[x] \)，使得
	\[ d = fu + gv. \]
\end{proposition}
\begin{proof}
\textbf{法一：}
若 \( \deg f < \deg g \)，由 \ref{带余除法}，存在 \( q_0, r_0 \in \mathbb{F}[x] \)，使得 \( f = q_0 g + r_0 \).
		
		那么有
		\[ (f, g) = (g, r_0)=d. \]
		
		重复这一过程：
		\[
		\begin{aligned}
			f &= q_0 g + r_0, \\
			g &= q_1 r_0 + r_1, \\
			&\vdots \\
			r_{n-3} &= q_{n-2} r_{n-2} + r_{n-1}, \\
			r_{n-2} &= q_{n-1} r_{n-1} + r_n, \\
			r_{n-1} &= q_n r_n.
		\end{aligned}
		\]
		\[(f,g)=(g,r_0)=(r_0,r_1)=\cdots=(r_{n-1},r_n)=cr_n\]
		其中 \( \deg r_n < \deg r_{n-1} < \cdots < \deg g <\deg f\)，所以存在 \( n \) 使得 \( r_{n-1}\mid r_n\) ，那么 \( d = r_n \).

		\[
		\begin{aligned}
			d &= r_{n-2} - q_{n-1} r_{n-1} \\
			&= r_{n-2} - q_{n-1}(r_{n-3} - q_{n-2} r_{n-2}) \\
			&\quad \vdots \\
			&= uf + vg.
		\end{aligned}
		\]
		
		这种证法给出了$d$的一种求法

\textbf{法二：} 考虑集合 \( \mathfrak{X} = \{ h(x) := m \cdot f + n \cdot g : m, n \in \mathbb{F}[x] \} \).

取$\mathfrak{X}$中次数最低的首一的多项式为 \( d(x) \).

下证：
1. \( d \mid f \) 且 \( d \mid g \)

2. 设 \( d_1 \mid f \) 且 \( d_1 \mid g \)，则 \( d_1 \mid d \)


1,\( f = q \cdot d + r \)，若 \( r \neq 0 \)，则 \( \deg r < \deg d \)，且 \( r = f - dq \in \mathfrak{X} \)，与 \( d \) 的取法矛盾.
故 \( r = 0 \)，即 \( d \mid f \).同理 \( d \mid g \).

2, 由于 \( d = uf + vg = vk_1 d_1 + uk_2 d_1 \Rightarrow d_1 \mid d \).

故 \( d = (f, g) \)存在且唯一.
\end{proof}
\begin{corollary}
	设 \( d = (f, g) \)，当且仅当 \( d \) 是 \( f, g \) 次数最大的公因式.
\end{corollary}
\begin{proof}
	必要性：设 \( d_1 \) 是 \( f, g \) 的公因式，且\( d_1 \mid d \)，那么 \( \deg d_1 \leqslant \deg d \).故 \( d \) 次数最大.
	
	充分性：设 \( \hat{d} = (f, g) \)，下证 \( d = \hat{d} \).
	
	由\ref{多项式环上的欧几里得算法}，存在 \( u, v \in \mathbb{F}[x] \)，使得 \( \hat{d} = uf + vg \).
	因为 \( d \) 是 \( f, g \) 的公因式，所以 \( d \mid \hat{d} \).
	
	从而 \( \deg d \leqslant \deg \hat{d} \leqslant \deg d \)，故 \( \deg d = \deg \hat{d} \)，故 \( d = \hat{d} \).
	
\end{proof}

欧几里得算法不依赖于数域的扩域.
\begin{corollary}
若 $f(x)$ 与 $g(x)$ 互素且 $f|h, g| h$ ，那么 $f g \mid h$
\end{corollary}
\begin{proof}

存在 \( p, q \in \mathbb{F}[x] \)，使得
\[ h = fp, \quad h = gq. \]

由 \( (f, g) = 1 \)，存在 \( u, v \in \mathbb{F}[x] \)，使得
\[ uf + vg = 1. \]

因此，
\[ ufh + vgh = h, \]
即
\[ ufgq + vgfp = h \Leftrightarrow fg(qu + pv) = h. \]

因此$f g \mid h$

\end{proof}

\begin{definition}\phantomsection\label{p1:concept:c1-l789}
	对于 \( f \in \mathbb{F}[x] \) 且 \( f \) 非零多项式，\( c \) 与 \( cf \) 都是 \( f(x) \) 的因子（\( c \neq 0 \) 且 \( c \in \mathbb{F} \)）.它们称为 \( f \) 的平凡因子.
	
	若一个多项式 \( p(x) \) 的因子只有平凡因子，则称 \( p(x) \) 为不可约的.相反的我们可以定义出可约的多项式
\end{definition}
多项式的可约性高度依赖于数域的选取，如$x^2-2$在$\Q$上不可约，但在$\R$上可约.

一条显而易见的结论是

\begin{proposition}
	设 \( p \in \mathbb{F}[x] \) 且是一个首一的不可约多项式.对于任意 \( f \in \mathbb{F}[x] \)，有
	\[
	(c p, f) =
	\begin{cases}
		p & \text{如果 } p \mid f, \\
		1 & \text{其他情况.}
	\end{cases}
	\]
\end{proposition}
\begin{proposition}
	若 \( f, g \in \mathbb{F}[x] \)，\( p \in \mathbb{F}[x] \) 且不可约，$p\mid fg$,那么 \( p \nmid f \) 或者 \( p \nmid g \).
	

\end{proposition}
\textbf{注意：} 逻辑学里的或者不是“要么A,要么B”的意思，而是“要么A,要么B,要么A且B”的意思
\begin{proof}
	设 \( p \nmid f \)，且 \( fg = p h \)，那么 \( (p, f) = 1 \).
	
	存在 \( u, v \in \mathbb{F}[x] \)，使得
	\[ fu + pv = 1. \]
	
	同乘
	\[ fgu + pgv = g. \]
	\[p(qu+gv)=g\]
	因此$p\mid g$
\end{proof}
由欧几里得算法我们得到以下结论
\begin{corollary}
	若 \( f_i \in \mathbb{F}[x] \) 不全为 0，\( f_1, \ldots, f_s \) 互素当且仅当存在 \( u_i \in \mathbb{F}[x] \) 使得
	\[ \sum_{i=1}^s u_i f_i = 1. \]
\end{corollary}

\begin{proof}
	\text{必要性：} 由\ref{多项式环上的欧几里得算法}，使用归纳法即可
	
	\text{充分性：}
	对于任意 \( i \)，\( d \mid f_i \) 那么 \( d \mid \sum u_i f_i = 1 \)，故 \( d = 1 \).
\end{proof}

\begin{corollary}
	\( f, g, h \in \mathbb{F}[x] \)，且 \( (g, f) = 1 \)，若 \( f \mid gh \)，则 \( f \mid h \).
	
\end{corollary}
\begin{proof}
	存在 \( u, v \in \mathbb{F}[x] \)，使得 \( gu + fv = 1 \)，设 \( fk = gh \).
	
	两边乘 \( h \)：
	\[ hgu + hfv = h, \]
	即
	\[ f(ku + hv) = h. \]
	
	故 \( f \mid h \).
	
	
\end{proof}
\begin{corollary}\label{1.3.14}

	若 \( (f, g_1) = 1 \)，\( (f, g_2) = 1 \)，则 \( (f, g_1 g_2) = 1 \).
	

\end{corollary}
\begin{proof}
	存在 \( u_1, v_1, u_2, v_2 \in \mathbb{F}[x] \)，使得
	\[ u_1 f + v_1 g_1 = 1, \]
	\[ u_2 f + v_2 g_2 = 1. \]
	
	相乘整理得：
	\[ (u_1 u_2 f + u_1 v_2 g_2 + u_2 v_1 g_1)f + (v_1 v_2) g_1 g_2 = 1. \]
	
	
\end{proof}
更一般地，
\[ (f, g) = 1 \Rightarrow (f^{n_1}, g^{n_2}) = 1. \]
\begin{corollary}
	若 \( (f_1, f_2) = 1 \)，且 \( f_1 \mid g \)，\( f_2 \mid g \)，则 \( f_1 f_2 \mid g \).
	
	
\end{corollary}
\begin{proof}
$$\begin{aligned}
	& f_1 u+f_2 v=1 \\
	& f_1 g u+f_2 g v=g \\
	& f_1 k_2 f_2 u+f_2 k_1 f_1 v=g \\
	& f_1 f_2 | g
\end{aligned}$$

\end{proof}
在多项式环上，我们也有因式分解定理
\begin{theorem}(因式分解定理)\label{因式分解定理}
对于任意一个首一多项式 \( f(x) \in \mathbb{F}[x] \)，\( f(x) \) 且非零多项式，则 \( f(x) \) 可以分解为有限个不可约多项式的乘积，即
	\[ f(x) = p_1 \cdots p_r = q_1 \cdots q_s, \]
	其中 \( p_i, q_j \) 是首一的不可约多项式.那么 \( r = s \)，且重排下有
	\[ p_1 = q_1, \ldots, p_r = q_s. \]
	
	

\end{theorem}
\begin{proof}
	存在性：对 \( \deg f \) 进行归纳.
	
	若 \( \deg f = 1 \)，\( f \) 已是一个不可约多项式，结论成立.
	
	假设对 \(\forall g\in\F[x] ,\deg g < \deg f \)，\( g \) 的因式分解存在.
	
	若 \( f \) 不可约，\(f=f\)便是一个分解.
	
	若 \( f \) 可约，设 \( f = gh \)，则 \( \deg g, \deg h < \deg f \)，且 \( g, h \) 可因式分解.则 \( f = gh \) 可因式分解.
	
	唯一性： 由 \( f = p_1 \cdots p_r = q_1 \cdots q_s \)，\( p_1 \mid q_1 \cdots q_s \)，存在$j$使$p_1\mid q_j$,重排使\(j=1\),同理\(q_1\mid p_1\) 故\( q_1 = p_1 \)，重复上述过程即可，故
	
	\[ p_1 = q_1, \ldots, p_r = q_r, \quad s = t. \]
	
	
	
\end{proof}
$\forall f$ 把上面的分解写成
\[f=ap_{1}^{r_1}\cdots p_{s}^{r_s}\]
\begin{definition}\phantomsection\label{p1:concept:c1-l916}
	对于 \( f, g \in \mathbb{F}[x] \)，我们称 \( g \) 是 \( f \) 的重因子，若 \( g^2 \mid f \)，若\( g^r \mid f \)且\( g^{r+1} \mid f \),则称 \( g \) 是 \( f \) 带有 \( r \) 重数的重因子.
\end{definition}

现在，我们考虑 \( f = a p_1^{r_1} \cdots p_s^{r_s} \)，\ref{因式分解定理}只给出了存在性，没有给出算法，下面给出一种算法，可以求出
\[ g = a p_1 \cdots p_s. \]



我们定义 \( \mathbb{F}[x] \) 上的导数算子
\begin{definition}\phantomsection\label{p1:concept:c1-l926}
	对于 \( f \in \mathbb{F}[x] \)，
	\[ f(x) = a_n x^n + a_{n-1} x^{n-1} + \cdots + a_1 x + a_0 \in \mathbb{F}[x], \]
	\( f(x) \) 的导数定义为
	\[ Df = f'(x) = n a_n x^{n-1} + (n-1) a_{n-1} x^{n-2} + \cdots + a_1 = \sum_{i=1}^n i a_i x^{i-1}. \]
	
	定义高阶导数\[ (f(x))' = (f'(x))' \]
	\[ (f^{(n)}(x))' = f^{(n+1)}(x) \]
\end{definition}
容易得到
\[ D(f + g) = Df + Dg \]
\[ D(fg) = f \cdot Dg + Df \cdot g \]

对于 \( n \in \mathbb{Z}_+ \)，有
\[ (f(x)^n)' = n f(x)^{n-1} f'(x). \]


\begin{proposition}
	\( p, f \in \mathbb{F}[x] \)，\( p \) 不可约，且
	\[ f(x) = p(x)^r h(x), \quad (p, h) = 1, \quad r \geq 1. \]
	
	那么
	\[ p^{r-1} \mid (f, f') \quad \text{且} \quad p^r \nmid (f, f'). \]
	
\end{proposition}
\begin{proof}
	\[ f' = r p^{r-1} p' h + p^r h' = p^{r-1} (r p' h + p h'). \]
	
	故
	\[ p^{r-1} \mid (f, f'). \]
	
	同时，由于 \( (p, h) = 1 \)，且$p$不可约
	\[(p,p')=1\]
	故
	\[ p\nmid r p' h + p h' \]
	因此
	\[ p^r \nmid (f, f'). \]
	

\end{proof}
\begin{corollary}\label{1.3.20}
	一个多项式 \( f \) 有重因子当且仅当
	\[ (f, f') \neq 1. \]
	
\end{corollary}
\begin{proof}
	设\( p \) 是 \( f \) 的不可约因子，那么
	\[ p \mid (f, f'). \]
	
	当且仅当
	\[ p^2 \mid f. \]
	
	由重因子的定义，得证.
	
\end{proof}

对于多项式 \( f = a p_1^{r_1} \cdots p_s^{r_s} \)，对于所有 \( i \)，有
\[ p_i^{r_i - 1} \mid f', \quad p_i^{r_i} \nmid f', \]
因此
\[ (f, f') = p_1^{r_1 - 1} \cdots p_s^{r_s - 1} \quad (r_i = 1 \text{ 时，} p_i \text{ 为平方因子}). \]

那么
\[ g = \frac{f}{(f, f')}, \]
则 \( g \) 为无重因子的多项式.
\section{多项式函数与根}
\begin{definition}\phantomsection\label{p1:concept:c1-l991}
	设 \[ f = a_n x^n + a_{n-1} x^{n-1} + \cdots + a_0 \in \mathbb{F}[x]\]，设 \(\mathbb{K}\text{为} \mathbb{F} \) 的任意扩域，\( \forall c \in \mathbb{K} \)， 有 \( \sum\limits_{i=0}^n a_i c^i \in \mathbb{K} \).
	
	从而有映射 \( f: \mathbb{K} \mapsto \mathbb{K}: c \mapsto \sum\limits_{i=0}^n a_i c^i \).
	
	称该映射为由 \( f \) 诱导的函数，仍记作 \( f(x) \).
	
	对于 \( f \in \mathbb{F}[x] \)，若 \( c \in \mathbb{C} \) 使得 \( f(c) = 0 \)，则称 \( c \) 为 \( f(x) \) 的一个根.
	
	对于 \( \alpha \in \mathbb{C} \)，若 \( \alpha \) 是整系数多项式的根，则称 \( \alpha \) 是代数数；否则，称 \( \alpha \) 是超越数.
\end{definition}


\begin{proposition}
	\( p, f \in \mathbb{F}[x] \)，\( p \) 在 \(\mathbb{F}\) 上不可约，若 \( p \text{与} f \) 有公共根，则 \( p \mid f \).
	
\end{proposition}
\begin{proof}
	 若\( p \nmid f \)，即 \( (p, f) = 1 \), \( \exists u, v \in \mathbb{F}[x] \)，使得 \( up + vf = 1 \).设 \( \alpha \) 为公共根，则 \( u(\alpha) p(\alpha) + v(\alpha) f(\alpha) = 0 \)，矛盾.
\end{proof}
\begin{proposition}
	\( (x - c) \mid f(x) \) 当且仅当 \( f(c) = 0 \).
\end{proposition}
\begin{proof}
	必要性：\( f(x) = g(x)(x - c) \)，则 \( f(c) = g(c)(c - c) = 0 \).
	
	充分性：\( f(x) = g(x)(x - c) + r \)，\( r \in \mathbb{F} \)，
	\[ f(c) = g(c)(c - c) + r \Rightarrow r = 0. \]
	

\end{proof}
若\( x - c \) 为 \( f(x) \) 的 \( k \) 重因式，即
\[ (x - c)^k \mid f(x), \quad (x - c)^{k+1} \nmid f(x). \]
\phantomsection\label{p1:concept:c1-l1024}则称 \( c \) 为 \( f(x) \) 的 \( k \) 重根.\( k = 1 \) 时，\( c \) 为单根；\( k > 1 \) 时，\( c \) 为重根.

一些显而易见的事实是：

1. \( f \) 在 \(\mathbb{F}\) 上有根是 \( f \) 在 \(\mathbb{F}\) 上可约的必要条件.

2. 当 \( \deg f = 2, 3 \) 时，\( f \) 在 \(\mathbb{F}\) 上可约可以推出 \( f \) 在 \(\mathbb{F}\) 上有根.

3. \( c \in \mathbb{C} \) 是重根，当且仅当 \( f(c) = f'(c) = 0 \).

\begin{lemma}
	对于 \( f \in \mathbb{F}[x] \)，若 \( \deg f = n \neq 0 \)，且 \( f \) 至多有 \( n \) 个根（计重数）.
	
\end{lemma}
\begin{proof}
	设 \( c_1, \ldots, c_s \) 是 \( f \) 生成互不相同的根，重数为 \( r_1, \ldots, r_s \).
	
	断言\((x-c_i)^{r_i}\mid f\),而
	\[ ((x - c_1)^{r_1}, \ldots, (x - c_s)^{r_s}) = 1, \]
	由\ref{1.3.14}，
	\[ (x - c_1)^{r_1} \cdots (x - c_s)^{r_s} \mid f. \]
	
	存在 \( g \in \mathbb{F}[x] \)，使得
	\[ g (x - c_1)^{r_1} \cdots (x - c_s)^{r_s} = f. \]
	
	因此，
	\[ \deg g +\deg r_1 + \cdots +\deg r_s = \deg f = n. \]
	
	所以，
	\[ r_1 + \cdots + r_s \leqslant n.\]
\end{proof}
\begin{corollary}
	\( f, g \in \mathbb{F}[x] \)，且 \( \deg f, \deg g < n \).若存在 \( n+1 \) 个互不相同的复数 \( c_1, \ldots, c_{n+1} \)，使得 \( f(c_i) = g(c_i) \)（\( 1 \leq i \leq n+1 \)），则 \( f = g \).
\end{corollary}
\begin{proof}
	 假设$f\neq g$,设 \( h = f - g \)，则 \( h \neq 0 \)，且 \( \deg h \leq \max(\deg f, \deg g) < n \).
	
	从而 \( h \) 至多有 \( n \) 个根（重根按重数计）.
	
	但 \( h(c_i) = f(c_i) - g(c_i) \)（\( 1 \leq i \leq n+1 \)），即 \( h \) 至少有 \( n+1 \) 个根，矛盾.因此，\( f = g \).
\end{proof}
\begin{corollary}
	\( f, g \in \mathbb{F}[x] \)，多项式 \( f = g \) 当且仅当即函数 \( f(x) = g(x) \)
\end{corollary}
\begin{proof}
	必要性：显然
	
	充分性：假设 \( f \neq g \)（作为多项式），设 \( h = f - g \)，则 \( h \neq 0 \).
	
	设 \( \deg h = n \).则 \( n \) 至少有 \( n \) 个根（重根按重数计）.
	
	但 \( h = 0 \)（作为函数），即 \( \forall c \in \mathbb{F} \)，都是多项式 \( h \) 的根，矛盾.所以 \( f = g \)（作为多项式）.
\end{proof}
下面我们介绍及其重要的代数基本定理，他又三种等价的形式，去证明设计复变函数，此处按下不表
\begin{theorem}(代数基本定理)\label{代数基本定理}
	任意一个$n$阶的多项式必有$n$个复根(重根按重数计)
\end{theorem}
\phantomsection\label{p1:concept:c1-l1081}我们称一个满足代数基本定理的域为代数闭域.
\begin{proposition}
	以下三条命题等价：
	\begin{enumerate}
		\item 任意 \( n \) 阶多项式恰有 \( n \) 个根（重根按重数计）.
		\item 任意多项式都有根.
		\item \( \mathbb{C} \) 上的不可约多项式次数为 1.
	\end{enumerate}
	
	
\end{proposition}
\begin{proof}
	我们证明 \( 1 \Rightarrow 2 \Rightarrow 1 \Rightarrow 3 \Rightarrow 1 \).
	
	\begin{itemize}
		\item \( 1 \Rightarrow 2 \)：显然成立.
		
		\item \( 2 \Rightarrow 1 \)：设 \( \deg f = n \)，由 2，设 \( f \) 有根 \( c \)，则 \( f(x) = g(x)(x - c) \)，且 \( \deg g = n - 1 \).对 \( \deg f \) 用归纳可证.
		
		\item \( 1 \Rightarrow 3 \)：设 \( \deg f = n \)，\( f \) 的根为 \( c_1, \ldots, c_n \)（可能有相同），则 \( f = a(x - c_1) \cdots (x - c_n) \).若 \( f \) 有不可约因式，则它必须是一次因式，否则矛盾.
		
		\item \( 3 \Rightarrow 1 \)：设 \( \deg f = n \)，\( f = p_1 \cdots p_n \)，且 \( \deg p_i = 1 \).则 \( \deg f = n \)，且 \( f \) 有 \( n \) 个根（重根按重数计）.
	\end{itemize}
\end{proof}

显而易见的，
多项式 \( f \) 不可约当且仅当 \( \deg f = 1 \)

\( f \) 在 \( \mathbb{C} \) 上的标准分解式都具备以下的形式：
\[ f(x) = a(x - c_1)^{r_1} \cdots (x - c_n)^{r_n}. \]

同时，给出 \( n \) 次单位根：
\[ x^n = 1. \]

则 \( x = \omega_k \)（\( 1 \leq k \leq n \)），其中
\[ \omega_k = e^{i\frac{ 2 k \pi}{n}}, \quad 1 \leqslant k \leqslant n.\]

\begin{corollary}
对于 \( f \in \mathbb{F}[x] \)，\( f \) 有重根当且仅当 \( (f, f') \neq 1 \).
\end{corollary}
\begin{proof}
	充分性：若 \( (f, f') = 1 \)，由 \ref{1.3.20}知 \( f \) 无重因式即无重根.
	
	必要性：若 \( (f, f') \neq 1 \)，存在 \( d \in \mathbb{F}[x] \) 为 \( f \) 的重因式，由 \ref{代数基本定理} 知，存在 \( c \in \mathbb{C} \) 使得 \( d(c) = 0 \)，则 \( c \) 是 \( f \) 的重根,矛盾.
	
\end{proof}
\section{$\C,\R,\Q$上的因式分解与不可约多项式}
在上一节中， 利用代数基本定理 ，我们得到了$\C$上的不可约多项式只有一次.下面,我们研究$\R,\Q$上的不可约多项式
\begin{proposition}
	实数域 $\mathbb{R}$ 上的不可约多项式的次数为 1 或 2.
\end{proposition}
\begin{proof}
	对于任意 $ p \in \mathbb{R}[x] $ 且 $ \deg p \geqslant 3 $，设 $ p $ 在 $\mathbb{R}$ 上不可约，则 $ p(x) $ 无实根.从而 $ p(x) $ 的根都是复根.
	
	对于 $ \alpha \in \mathbb{C}\setminus\R $，若 $ p(\alpha) = 0 $，则
	$$
	p(\overline{\alpha}) = \overline{p(\alpha)} = 0.
	$$
	
	因此，$ \overline{\alpha }$ 也是 $ p $ 的根.
	
	$ p(x) $ 有因式 $ (x - \alpha)(x - \overline{\alpha}) = x^2 - (\alpha + \overline{\alpha})x + \alpha \overline{\alpha} \in \mathbb{R}[x] $，
	因此 $ p(x) $ 在 $\mathbb{R}$ 上可约.矛盾
	
\end{proof}

综上，我们得出 $\mathbb{R}[x]$ 上的标准分解式：
$$
f(x) = a (x - a_1)^{r_1} \cdots (x - a_s)^{r_s} (x^2 + p_1 x + q_1)^{k_1} \cdots (x^2 + p_t x + q_t)^{k_t},
$$
其中，$ a_i \in \mathbb{R} $ 且互不相同，$ (p_j, q_j) \in \mathbb{R}^2 $ 且不相同，且 $ p_j^2 - 4 q_j < 0 $，$ 1 \leq j \leq t $.

下面我们考虑$\Q$上的多项式，对于 \( \mathbb{Q}[x] \) 中的多项式 \( f \)，总有整系数多项式与之相伴，考虑整系数多项式即可.
\begin{definition}\phantomsection\label{p1:concept:c1-l1154}
	$$ f(x) = a_n x^n + \cdots + a_0 \in \mathbb{Z}[x]$$\( \)\( a_n \neq 0 \),若 \( (a_n, \ldots, a_0) = 1 \)，称 \( f \) 为本原多项式.
\end{definition}

\begin{lemma}\label{1.5.3}
	$\forall f \in \mathbb{Q}[x], \exists!g(x)$ 与 $f(x)$ 相伴且为本原多项式
\end{lemma}
\begin{proof}
	存在性显然，下证明唯一性，也即证明若两个本原多项式相伴，则 $g_1= \pm g_2$
	
	设 \( g_1 = \sum b_i x^i \)，\( g_2 = \sum c_i x^i \)，存在 \( \dfrac{s}{r} \neq 0 \) 且 \( \dfrac{s}{r} \in \mathbb{Q} \) 且 \( (s, r) = 1 \)，使得
	\[ g_1 = \frac{s}{r} g_2, \]
	即 \( b_i = \dfrac{s}{r} c_i \)，对于 \( 0 \leqslant i \leqslant n \) 也成立，即 \( r b_i = s c_i \).
	
	可以推出 \( r \mid s c_i \) 而\( (r, s) = 1 \)，故 \( r \mid c_i \)，即 \( r \mid (c_1, \ldots, c_n) = 1 \).
	
	因此，\( r = \pm 1 \)，同理 \( s = \pm 1 \)，故 \( g_1 = \pm g_2 \).
\end{proof}
\begin{theorem}(高斯引理)\label{高斯引理}
	若 \( f, g \) 是整系数本原多项式，则 \( h = fg \) 也是整系数本原多项式.
\end{theorem}
\begin{proof}
	设 \( f = \sum\limits_{i=0}^n a_i x^i \)，\( g = \sum\limits_{j=0}^m b_j x^j \) 是整系数本原多项式.
	
	任取 素数\( p \) ，下证 \( p \) 不是 \( fg \) 系数的公因子，这等价于系数中至少有一个无法被 \( p \) 整除.
	
	由于 \( f \) 是本原多项式，所以 \( p \) 不整除 \( a_0, \ldots, a_n \) 的公因子.从而存在 \( i \)，使得
	\[ p \mid a_0, p \mid a_1,\cdots, p \mid a_{i-1}\text{但} p \nmid a_i,0\leqslant i\leqslant n \]
	
	对于 \( g \)，存在 \( j \)，使得
	\[ p \mid b_0, p \mid b_1,\cdots, p \mid b_{j-1}\text{但} p \nmid b_j,0\leqslant j\leqslant n \]
	
	现在考虑 \( fg \) 的系数：
	\[ fg = \sum_{k=0}^{n+m} \left( \sum_{r+s=k} a_r b_s \right) x^k. \]
	
	考虑 \( x^{i+j} \) 的系数：
	\[ c_{i+j} = a_i b_j + a_{i-1} b_{j+1} + \cdots + a_0 b_{i+j} + \cdots + a_{i+j} b_0. \]
	
	由于 \( p \nmid a_i \) 和 \( p \nmid b_j \)，所以 \( p \nmid a_i b_j \)，即 \( p \) 不整除 \( fg \) 的系数$c_{i+j}$.因此，\( fg \) 是本原多项式.
	
\end{proof}
\begin{corollary}
	若 \( f, g \) 是本原多项式，且 \( \exists h(f = gh) \)，则 \( h \) 是整系数多项式。
\end{corollary}

\begin{proof}
	由于 \( h \in \mathbb{Q}[x] \)，存在本原多项式 \( h_1 \) 与 \( h \) 相伴。由 \( f = gh = c g h_1 \)，
	根据\ref{高斯引理}，\( g h_1 \) 也是本原多项式。因此 \( g h_1 \) 与 \( f \) 相伴。由\ref{1.5.3}知，\( c = \pm 1 \)，故 \( h = \pm h_1 \)，
	即 \( h \in \mathbb{Z}[x] \)。
\end{proof}
也就是说，所有本原多项式构成的集合是一个乘法交换半群。
\begin{corollary}
	若 \( f \in \mathbb{Z}[x] \) 在 \(\mathbb{Q}\) 上可约，则 \( f \) 可以分解为次数更低的整系数多项式的乘积，等价地说， \(\mathbb{Q}\) 上的可约等价于在 \(\mathbb{Z}\) 上的可约。
\end{corollary}
\begin{proof}
	设 \( f \in \mathbb{Z}[x] \) 且 \( f \) 在 \(\mathbb{Q}\) 上可约。那么存在 \( g, h \in \mathbb{Q}[x] \)，使得 \( f = gh \)。 \( \deg g, \deg h < \deg f \)
	
	存在 \( c, c_1, c_2 \in \mathbb{Q} \) 且 \( c_1, c_2 \neq 0 \)，使得 \( f = c f_1 \), \( g = c_1 g_1 \), \( h = c_2 h_1 \)，其中 \( f_1, g_1, h_1 \) 是本原多项式。
	
	根据 \( c f_1 = c_1 c_2 g_1 h_1 \)，由\ref{1.5.3}及\ref{高斯引理}，有 \( c = \pm c_1 c_2 \)，因此
	\[ f = gh = (\pm  g_1) h_1, \]
	其中 \( h_1, \pm  g_1 \in \mathbb{Z}[x] \)，且
	\[ \deg (\pm g_1) < \deg f,  \deg h_1 < \deg f. \]
	
\end{proof}
\begin{proposition}
	设 \( f(x) = a_n x^n + \dots + a_1 x + a_0 \) 是一个整系数多项式，
	若 \( x = \dfrac{s}{r} \)（其中 \( s,r \in \mathbb{Z} \)，且 \( (s,r)=1 \)）是 \( f(x) \) 的一个有理根，则
	\[
	s \mid a_0,\quad r \mid a_n
	\]
\end{proposition}
\begin{proof}
	若 \( \dfrac{s}{r} \) 是 \( f(x) \) 的一个根，则
	\[
	\left(x-\frac{s}{r}\right) \mid f(x) \iff (rx-s) \mid f(x)
	\]
	因此存在多项式 \( h(x) \in \mathbb{Z}[x] \)，使得
	\[
	(rx-s)h(x) = f(x)
	\]
	设 \( h(x) = \sum\limits_{j=0}^{n-1} b_j x^j \)，比较两边多项式的系数可得：
	\[
	\begin{cases}
		r b_{n-1} = a_n \\
		-s b_0 = a_0
	\end{cases}
	\]
	由此可知 \( r \mid a_n \)，且 \( s \mid a_0 \)。
	
\end{proof}
\begin{theorem}（Eisenstein 判别法）
	设
	\[
	f(x)=a_nx^n+\dots+a_0\in\mathbb{Z}[x]
	\]
	若存在素数 \(p\)，使得
	\begin{enumerate}
		\item \(p\mid a_0\)，但 \(p^2\nmid a_0\)；
		\item \(p\mid a_i\)，对 \(1\le i\le n-1\)；
		\item \(p\nmid a_n\)。
	\end{enumerate}
	那么 \(f(x)\) 在 \(\mathbb{Q}\) 上不可约。
\end{theorem}
\begin{proof}
	反证法。若不然，\(f(x)\) 在 \(\mathbb{Q}\) 上可约，存在 \(g,h\in\mathbb{Z}[x]\)，使得 \(f=gh\)，且 \(\deg g<\deg f\)，\(\deg h<\deg f\)。
	
	设
	\[
	g=\sum_{j=0}^m b_jx^j,\quad h=\sum_{k=0}^l c_kx^k,\quad m<n,\ l<n
	\]
	
	由 \(a_0=b_0c_0\) 且 \(p\mid a_0,\ p^2\nmid a_0\)，可知 \(p\) 恰好整除 \(b_0,c_0\) 中的一个，不妨设 \(p\mid b_0,\ p\nmid c_0\)。
	
	由 \(a_n=b_mc_l\) 且 \(p\nmid a_n\)，可知 \(p\nmid b_m,\ p\nmid c_l\)。
	
	因为 \(p\mid b_0\)，\(p\nmid b_m\)，所以存在 \(i\)（\(1\le i\le m\)），使得 \(p\mid b_0,p\mid b_1,\dots,p\mid b_{i-1}\) 但 \(p\nmid b_i\)。
	
	考虑 \(a_i=b_ic_0+b_{i-1}c_1+\dots\)，右边除了 \(b_ic_0\) 这一项外，其余项都能被 \(p\) 整除，而 \(p\nmid b_ic_0\)，故 \(p\nmid a_i\)，这与条件中 \(p\mid a_i\)（\(1\le i\le n-1\)）矛盾。
	
	因此，\(f(x)\) 在 \(\mathbb{Q}\) 上不可约。
	
\end{proof}
\section{多元多项式环}
\begin{definition}\phantomsection\label{p1:concept:c1-l1278}
	给定一个数域 \(\F \)，\( n \in \mathbb{Z}^+, n>1 \)，取不定元 \( x_1,\dots,x_n \)。
	
	定义\textbf{单项式}是形如 \( a x_1^{i_1} \cdots x_n^{i_n} \) 的表达式，其中 \( i_j \) 是非负整数，\( x_i \notin \F \)，\( a \neq 0 \)。
	\begin{enumerate}
		\item \( a \) 称为该单项式的\textbf{系数}。
		\item \( n \) 元组 \( (i_1,\dots,i_n) \) 称为单项式的\textbf{指数组}。
		\item \( i_1+\dots+i_n \) 称为该单项式的\textbf{次数}。
	\end{enumerate}
	
	
	规定：\( a x_1^{i_1} \cdots x_n^{i_n} = b x_1^{j_1} \cdots x_n^{j_n} \) 当且仅当 \( a=b \) 且 \( i_k=j_k \)（\( 1 \le k \le n \)）。
	
	同类项规定为：指数组相同（\( i_k=j_k \)）的项，系数 \( a,b \) 可以不相等。
	
	称有限和
	\[
	f(x_1,\dots,x_n)=\sum_{i_1,\dots,i_n} a_{i_1,\dots,i_n} x_1^{i_1} \cdots x_n^{i_n}
	\]
	为一个\textbf{多元多项式}，零多项式记为 \( 0 \)，\( a_{i_1,\cdots,i_n} \) 为 \( f \) 的系数。
	
	规定
	\[
	\sum_{i_1,\dots,i_n} a_{i_1,\dots,i_n} x_1^{i_1} \cdots x_n^{i_n} = f = g = \sum_{j_1,\dots,j_n} b_{j_1,\dots,j_n} x_1^{j_1} \cdots x_n^{j_n}
	\]
	当且仅当 \( \forall (i_1,i_2,\dots,i_n) \)，都有 \( a_{i_1,\dots,i_n} = b_{i_1,\dots,i_n} \)。
	
	所有 \( n \) 元多项式构成一个集合 \( \F[x_1,\dots,x_n] \)，在我们之后规定的运算下，它是一个\textbf{多项式环}。
	
	若\[
	\sum_{i_1,\dots,i_n} a_{i_1,\dots,i_n} x_1^{i_1} \cdots x_n^{i_n} = 0 \iff \forall (i_1,\dots,i_n),\ a_{i_1,\dots,i_n}=0
	\]
	
	也就是说，\( x_1,\dots,x_n \) 在 \( F \) 上是\textbf{代数无关}的。代数无关是一个强于线性无关的概念，举例来说：集合 \( \{1,\sin x,\cos x\} \) 代数相关，但在实数域上线性无关。
	
\end{definition}
显然的，$\F[x_1,\cdots,x_{n-1}]\subset\F[x_1,\cdots,x_{n}]$
\begin{definition}\phantomsection\label{p1:concept:c1-l1315}
	当 \( f(x_1,\dots,x_n) \neq 0 \) 时，
	\[
	\deg f: = \max\{i_1+\dots+i_n \mid a_{i_1,\dots,i_n} \neq 0\}
	\]
	定义为 \( f \) 的次数，记为$\deg f$,也可记为 \( \partial f \)。
	规定零多项式的次数为 \( \deg 0 = -\infty \)。
\end{definition}
为了在 \( \F[x_1,\dots,x_n] \) 构建某些概念，我们在 \( \F[x_1,\dots,x_n] \) 上定义字典序。
\begin{definition}\phantomsection\label{p1:concept:c1-l1324}
	对指数组 \( (i_1,\dots,i_n) \) 和 \( (j_1,\dots,j_n) \)，若存在 \( k \)，使得 \( i_1=j_1,\dots,i_{k-1}=j_{k-1} \)，且 \( i_k>j_k \)，就称 \( (i_1,\dots,i_n) \) 大于 \( (j_1,\dots,j_n) \)，记为 \( (i_1,\dots,i_n) > (j_1,\dots,j_n) \)。
	
	这种序关系我们称为\textbf{字典序}。一个多项式的字典序最大的单项式称为\textbf{首项}。
	
\end{definition}


接下来，我们定义 \( \F[x_1,\dots,x_n] \) 上的二元运算。

\begin{definition}\phantomsection\label{p1:concept:c1-l1334}
	设
	\[
	f = \sum_{i_1,\dots,i_n} a_{i_1,\dots,i_n} x_1^{i_1} \cdots x_n^{i_n}, \quad g = \sum_{j_1,\dots,j_n} b_{j_1,\dots,j_n} x_1^{j_1} \cdots x_n^{j_n}
	\]
	则加法定义为：
	\[
	f+g := \sum_{k_1,\dots,k_n} (a_{k_1,\dots,k_n} + b_{k_1,\dots,k_n}) x_1^{k_1} \cdots x_n^{k_n}
	\]
	乘法定义为：
	\[
	\begin{aligned}
		f \cdot g &:= \left( \sum_{i_1,\dots,i_n} a_{i_1,\dots,i_n} x_1^{i_1} \cdots x_n^{i_n} \right) \left( \sum_{j_1,\dots,j_n} b_{j_1,\dots,j_n} x_1^{j_1} \cdots x_n^{j_n} \right) \\
		&= \sum_{k_1,\dots,k_n} \left( \sum_{\substack{i_1+j_1=k_1 \\ \dots \\ i_n+j_n=k_n}} a_{i_1,\dots,i_n} b_{j_1,\dots,j_n} \right) x_1^{k_1} \cdots x_n^{k_n}
	\end{aligned}
	\]
	显然地，\( \F[x_1,\dots,x_n] \)在上述的加法与乘法中，构成一个环。
\end{definition}

下面考虑 \( \F[x_1,\dots,x_n] \) 上是否有消去律。考虑首项，我们得到以下命题，它与消去律等价。
\begin{proposition}
	在 \( \F[x_1,x_2,\cdots,x_n] \) 中，两个多项式乘积的首项等于它们首项的乘积。
\end{proposition}
\begin{proof}
	设 \( \alpha = a_{i_1,\dots,i_n} x_1^{i_1} \cdots x_n^{i_n} \) 为 \( f \) 的首项，\( \beta = b_{j_1,\dots,j_n} x_1^{j_1} \dots x_n^{j_n} \) 为 \( g \) 的首项，则 \( a_{i_1,\dots,i_n} \neq 0 \)，\( b_{j_1,\dots,j_n} \neq 0 \)，从而 \( a_{i_1,\dots,i_n} b_{j_1,\dots,j_n} \neq 0 \)。
	
	要证 \( \alpha\beta = a_{i_1,\dots,i_n} b_{j_1,\dots,j_n} x_1^{i_1+j_1} \cdots x_n^{i_n+j_n} \) 是 \( fg \) 的首项，只需证明指数组 \( (i_1+j_1,\dots,i_n+j_n) \) 大于 \( fg \) 中其他任意指数组。
	
	分三种情况讨论：
	\begin{enumerate}
		\item \( (i_1+j_1,\dots,i_n+j_n) > (i_1+j'_1,\dots,i_n+j'_n) \)，其中 \( (j_1,\dots,j_n) > (j'_1,\dots,j'_n) \)；
		\item \( (i_1+j_1,\dots,i_n+j_n) > (i'_1+j_1,\dots,i'_n+j_n) \)，其中 \( (i_1,\dots,i_n) > (i'_1,\dots,i'_n) \)；
		\item \( (i_1+j_1,\dots,i_n+j_n) > (i'_1+j'_1,\dots,i'_n+j'_n) \)，其中 \( (i_1,\dots,i_n) > (i'_1,\dots,i'_n) \) 且 \( (j_1,\dots,j_n) > (j'_1,\dots,j'_n) \)。
	\end{enumerate}
	
	由字典序的定义，这三种情况都成立。
\end{proof}


注意：多项式的首项次数不一定是多项式的次数。例如，\( x_1 - x_2^{100} \) 的首项是 \( x_1 \)，次数为 1，但多项式的次数为 100。

为了定义乘积的次数，我们考虑齐次多项式。



\begin{definition}\phantomsection\label{p1:concept:c1-l1379}
	若多项式 \( f(x_1,\dots,x_n) \) 中每个单项式的次数都相同且等于 \( m \)，则称 \( f \) 是一个 \( m \) 次齐次多项式。
\end{definition}
例如
\begin{enumerate}
	\item \( f = \sum\limits_{i=1}^n a_i x_i \) 是一次齐次多项式；
	\item\phantomsection\label{p1:concept:c1-l1385}\( f = \sum\limits_{i=1}^n x_i^k \) 是 \( k \) 次齐次多项式，记为 \( S_k \)，称为 \textbf{Newton 和}。  
\end{enumerate}
\( m \) 次齐次多项式有以下性质：
\[
f(tx_1,\dots,tx_n) = t^m f(x_1,\dots,x_n)
\]

对任意 \( f \in F[x_1,\dots,x_n] \)，设 \( \deg f = s \)。令 \( f_i \) 表示 \( f \) 中所有次数为 \( i \) 的单项式组成的齐次多项式，则 \( f \) 可唯一表示为：
\[
f = f_0 + f_1 + \dots + f_s
\]
\begin{definition}
	对任意 \( f,g \in F[x_1,\dots,x_n] \)，多项式乘积的次数满足：
	\[
	\deg(fg) = \deg f + \deg g
	\]
\end{definition}
\begin{proof}
	设 \( \deg f = m \)，\( \deg g = n \)，则 \( f \) 和 \( g \) 可以分解为齐次多项式之和：
	\[
	f = f_0 + f_1 + \dots + f_m, \quad g = g_0 + g_1 + \dots + g_n
	\]
	其中 \( f_i \) 是 \( f \) 的 \( i \) 次齐次部分，\( g_i \) 是 \( g \) 的 \( i \) 次齐次部分，且 \( f_m \neq 0 \)，\( g_n \neq 0 \)。
	
	将它们相乘，得到：
	\[
	fg = f_m g_n + (f_{m-1}g_n + f_m g_{n-1}) + \dots + (f_0 g_1 + f_1 g_0) + f_0 g_0
	\]
	其中 \( f_m g_n \) 是次数为 \( m+n \) 的齐次多项式，且 \( f_m g_n \neq 0 \)。而其余项的次数均小于 \( m+n \)。
	
	因此，乘积 \( fg \) 的次数就是 \( m+n \)，即：
	\[
	\deg(fg) = \deg(f_m g_n) = m + n = \deg f + \deg g
	\]
\end{proof}
在 \( \F[x_1,\dots,x_n] \) 上，可以定义整除、因子、最大公因式、不可约多项式的概念。
例如，对 \( \mathbb{C}[x,y] \)，多项式 \( x^n - y \) 在 \( \mathbb{C} \) 上不可约。
多元多项式环同样成立唯一分解定理。

但是，\( \F[x_1,\dots,x_n] \) 上没有带余除法（当 \( n>1 \) 时），也就是说，\( \F[x_1,x_2,\dots,x_n] \) 不是一个欧几里得环，而是一个唯一分解环（UFD）。


类似地，我们定义多元多项式函数。
\begin{definition}\phantomsection\label{p1:concept:c1-l1428}
	设 \( \F \subseteq \mathbb{K} \)，\( \mathbb{K}\) 为 \( \F \) 的扩域。
	对多项式
	\[
	f(x_1,\dots,x_n) = \sum_{i_1,\dots,i_n} a_{i_1,\dots,i_n} x_1^{i_1} \cdots x_n^{i_n} \in \F[x_1,\dots,x_n]
	\]
	定义多项式函数：
	\[
	\begin{aligned}
		f: \mathbb{K}^n &\to \mathbb{K} \\
		(a_1,\dots,a_n) &\mapsto f(a_1,\dots,a_n)
	\end{aligned}
	\]
\end{definition}


显然地，以下两命题成立且等价：

1.  对 \( f,g \in \F[x_1,\dots,x_n] \)，多项式 \( f=g \) 当且仅当对应的函数 \( f=g \)。

2.  多项式为零多项式，当且仅当其对应的函数为零函数。

对任意 \( f \in \F[x_1,\dots,x_n] \)，可以写成以下形式：
\[
f(x_1,\dots,x_n) = f_0(x_1,\dots,x_{n-1}) + f_1(x_1,\dots,x_{n-1})x_n + \dots + f_m(x_1,\dots,x_{n-1})x_n^m
\]
其中 \( m \) 是 \( x_n \) 在 \( f \) 中的最高次数。
也就是说，多项式环可以写成：
\[
\F[x_1,\dots,x_n] = \F[x_1,\dots,x_{n-1}][x_n]
\]
\section{对称多项式}
\begin{definition}\phantomsection\label{p1:concept:c1-l1460}
	设 \( f(x_1,\dots,x_n) \F[x_1,\dots,x_n] \)，如果对任意 \( i,j \)，都有
	\[
	f(x_1,\dots,x_i,\dots,x_j,\dots,x_n) = f(x_1,\dots,x_j,\dots,x_i,\dots,x_n)
	\]
	则称 \( f \) 是一个 \( n \) 元对称多项式。
	
	定义 \( S[x_1,\dots,x_n] \subset \F[x_1,\dots,x_n] \) 为对称多项式的集合。
\end{definition}


考虑一元多项式 \( f(x)=a_nx^n+\dots+a_0 \in \F[x] \)，由代数基本定理，它在 \( \mathbb{C} \) 中有根 \( x_1,\dots,x_n \)。
那么
\[
\begin{aligned}
	f(x) &= a_n(x-x_1)\cdots(x-x_n) \\
	&= a_n\left[ x^n + \left(-\sum_{i=1}^n x_i\right)x^{n-1} + \left(\sum_{1\le i<j\le n} x_i x_j\right)x^{n-2} + \dots \right. \\
	& \left. + (-1)^k \sum_{1\le i_1<\dots<i_k\le n} \prod_{m=1}^k x_{i_m} x^{n-k} + \dots + (-1)^n x_1\cdots x_n \right]
\end{aligned}
\]
比较系数，得到
\[
\sum_{i=1}^n x_i = -\frac{a_{n-1}}{a_n}
,\sum_{1\le i<j\le n} x_i x_j = \frac{a_{n-2}}{a_n},
\dots,
\sum_{1\le i_1<\dots<i_k\le n} x_{i_1}\cdots x_{i_k} = (-1)^k\frac{a_{n-k}}{a_n},
\dots,
x_1\cdots x_n = (-1)^n\frac{a_0}{a_n}
\]
这就是著名的 Vieta 定理。

将 \( x_1,\dots,x_n \) 视作不定元，我们得到以下定义：

\begin{definition}\phantomsection\label{p1:concept:c1-l1493}
	令
	\[
	\begin{aligned}
		\sigma_1 &= \sum_{i=1}^n x_i, \\
		\sigma_2 &= \sum_{1\le i<j\le n} x_i x_j, \\
		&\dots \\
		\sigma_k &= \sum_{1\le i_1<\dots<i_k\le n} x_{i_1}\cdots x_{i_k}, \\
		&\dots \\
		\sigma_n &= x_1\cdots x_n.
	\end{aligned}
	\]
	\( \sigma_1,\dots,\sigma_n \) 都是对称多项式，我们称之为基本对称多项式。
\end{definition}

\begin{theorem}（对称多项式基本定理）
	对任意 \( f(x_1,\dots,x_n) \in \F[x_1,\dots,x_n] \)，有
	\[
	f \in S[x_1,\dots,x_n] \iff \exists! g \in \F[x_1,\dots,x_n]
	\]
	使得 \( g(\sigma_1,\dots,\sigma_n) = f(x_1,\dots,x_n) \)。
	
\end{theorem}

这条定理有以下两种等价形式：
\begin{enumerate}
	\item  \( \sigma_1,\dots,\sigma_n \) 代数独立；
	\item \( S[x_1,\dots,x_n] \) 是 \( \F[x_1,\dots,x_n] \) 的子环，且 \( S[x_1,\dots,x_n] = \F[\sigma_1,\dots,\sigma_n] \)，即 \( S \) 是 \( \sigma_1,\dots,\sigma_n \) 的多项式环。
\end{enumerate}


\begin{proof}
	\( \Leftarrow \)：显然成立，因为 \( g(\sigma_1,\dots,\sigma_n) \) 是对称多项式。
	
	\( \Rightarrow \)（存在性）：不妨先设 \( f \) 为 \( m \) 次齐次多项式（一般情况可分解为齐次多项式之和）。
	设 \( f \) 的首项为 \( a x_1^{i_1} \cdots x_n^{i_n} \)，由字典序定义，必有 \( i_1 \ge i_2 \ge \dots \ge i_n \) 且 \( i_1+\dots+i_n=m \)。
	
	考虑单项式 \( a \sigma_1^{t_1} \cdots \sigma_n^{t_n} \)，其首项为：
	\[
	a x_1^{t_1+\dots+t_n} \cdots x_n^{t_n}
	\]
	令其与 \( f \) 的首项相同，可得方程组：
	\[
	\begin{cases}
		t_1+t_2+\dots+t_n = i_1 \\
		t_2+\dots+t_n = i_2 \\
		\quad \vdots \\
		t_n = i_n
	\end{cases}
	\Rightarrow
	\begin{cases}
		t_1 = i_1 - i_2 \\
		t_2 = i_2 - i_3 \\
		\quad \vdots \\
		t_{n-1} = i_{n-1} - i_n \\
		t_n = i_n
	\end{cases}
	\]
	因 \( i_1 \ge i_2 \ge \dots \ge i_n \)，故 \( t_k \ge 0 \)
	
	构造多项式 \( f_1 = f - a \sigma_1^{i_1-i_2} \cdots \sigma_n^{i_n} \)，它是对称多项式，且其首项在字典序下严格小于 \( f \) 的首项。
	若 \( f_1=0 \)，则 \( f = a \sigma_1^{i_1-i_2} \cdots \sigma_n^{i_n} \)；
	若 \( f_1 \neq 0 \)，对 \( f_1 \) 重复上述操作。由于字典序下次数是有限的，经过有限步后，必然会得到零多项式。因此 \( f \) 可表示为基本对称多项式的多项式，存在性得证。


	(唯一性)：
	唯一性等价于：若 \( g(\sigma_1,\dots,\sigma_n)=0 \)，则 \( g=0 \in \F[x_1,\dots,x_n] \)。
	
	考虑多项式映射 \( \varphi: \mathbb{C}^n \to \mathbb{C}^n \)，其中 \( \varphi(x_1,\dots,x_n) = (\sigma_1,\dots,\sigma_n) \)。
	若 \( g(\sigma_1,\dots,\sigma_n)=0 \)，则复合函数 \( g \circ \varphi = 0 \)。要证 \( g=0 \)，只需证 \( \varphi \) 是满射。
	
	对任意 \( (z_1,\dots,z_n) \in \mathbb{C}^n \)，考虑多项式：
	\[
	h(x) = x^n - z_1 x^{n-1} + z_2 x^{n-2} - \dots + (-1)^n z_n
	\]
	由代数基本定理，\( h(x) \) 在 \( \mathbb{C} \) 上有根 \( \alpha_1,\dots,\alpha_n \)。根据 Vieta 定理，这些根满足：
	\[
	\begin{cases}
		\sigma_1(\alpha_1,\dots,\alpha_n) = z_1 \\
		\quad \vdots \\
		\sigma_n(\alpha_1,\dots,\alpha_n) = z_n
	\end{cases}
	\]
	因此 \( \varphi(\alpha_1,\dots,\alpha_n) = (z_1,\dots,z_n) \)，故 \( \varphi \) 是满射。因此 \( g \) 在 \( \mathbb{C}^n \) 上恒为零，从而 \( g=0 \)，唯一性得证。
	
\end{proof}
\begin{example}
	将多项式 \( (x_1^2 + x_2x_3)(x_2^2 + x_1x_3)(x_3^2 + x_1x_2) \) 表示为基本对称多项式 \( \sigma_1, \sigma_2, \sigma_3 \) 的多项式。
	
	
\end{example}
\begin{solution}
多项式的首项为 \( x_1^4x_2x_3 \)，对应的指数组为 \( (4,1,1) \)。
我们列出所有满足 \( j_1 \ge j_2 \ge j_3 \ge 0 \)、\( j_1+j_2+j_3=6 \) 且 \( (4,1,1) \ge (j_1,j_2,j_3) \) 的指数组，并写出对应的基本对称多项式单项式：
\[
\begin{aligned}
	(4,1,1) &\Rightarrow (3,0,1) \Rightarrow \sigma_1^3\sigma_3 \\
	(3,3,0) &\Rightarrow (0,3,0) \Rightarrow \sigma_2^3 \\
	(3,2,1) &\Rightarrow (1,1,1) \Rightarrow \sigma_1\sigma_2\sigma_3 \\
	(2,2,2) &\Rightarrow (0,0,2) \Rightarrow \sigma_3^2
\end{aligned}
\]
因此，该多项式可以表示为：
\[
f(x_1,x_2,x_3) = (x_1^2 + x_2x_3)(x_2^2 + x_1x_3)(x_3^2 + x_1x_2) = \sigma_1^3\sigma_3 + b\sigma_2^3 + c\sigma_1\sigma_2\sigma_3 + d\sigma_3^2
\]
接下来通过代入特殊值来确定系数 \(b, c, d\)：

取 \(x_1=x_2=1, x_3=0\)，则 \( \sigma_1=2, \sigma_2=1, \sigma_3=0, f=1 \)。代入得：
\(1 = 2^3 \cdot 0 + b \cdot 1^3 + c \cdot 2 \cdot 1 \cdot 0 + d \cdot 0^2 \Rightarrow b=1\)。

取 \(x_1=2, x_2=x_3=-1\)，则 \( \sigma_1=0, \sigma_2=-3, \sigma_3=2, f=5 \)。代入得：
\(5 = 0^3 \cdot 2 + 1 \cdot (-3)^3 + c \cdot 0 \cdot (-3) \cdot 2 + d \cdot 2^2 \Rightarrow 5 = -27 + 4d \Rightarrow d=8\)。

 取 \(x_1=x_2=x_3=1\)，则 \( \sigma_1=3, \sigma_2=3, \sigma_3=1, f=8 \)。代入得：
\(8 = 3^3 \cdot 1 + 1 \cdot 3^3 + c \cdot 3 \cdot 3 \cdot 1 + 8 \cdot 1^2 \Rightarrow 8 = 27 + 27 + 9c + 8 \Rightarrow c=-6\)。

综上，该多项式用基本对称多项式表示为：
\[
f(x_1,x_2,x_3) = \sigma_1^3\sigma_3 + \sigma_2^3 - 6\sigma_1\sigma_2\sigma_3 + 8\sigma_3^2
\]



\end{solution}
\section{一些应用：关于多项式}
\begin{proposition}（拉格朗日插值公式）
	设 \( a_1,\dots,a_n \) 是 \( n \) 个互不相同的数，定义多项式
	\[
	g(x) = (x-a_1)\cdots(x-a_n)
	\]
	对任意给定的 \( b_1,\dots,b_n \)，存在唯一的多项式 \( L(x) \in F[x] \)，满足 \( \deg L < n \)，且
	\[
	L(a_i) = b_i,\quad i=1,2,\dots,n
	\]
	该多项式由以下公式给出：
	\[
	L(x) = \sum_{i=1}^n b_i \prod_{\substack{k=1 \\ k\neq i}}^n \frac{x-a_k}{a_i-a_k} = \sum_{i=1}^n \frac{b_i g(x)}{(x-a_i)g'(a_i)}
	\]
	
\end{proposition}

\begin{proof}
	首先，证明一个恒等式：
	\[
	\sum_{i=1}^n \frac{g(x)}{(x-a_i)g'(a_i)} = 1
	\]
	设 \( F(x) = \sum\limits_{i=1}^n \dfrac{g(x)}{(x-a_i)g'(a_i)} - 1 \)。
	由 \( g(x) = \prod\limits_{k=1}^n (x-a_k) \)，其导数为 \( g'(x) = \sum\limits_{i=1}^n \prod\limits_{\substack{k=1 \\ k\neq i}}^n (x-a_k) \)，故 \( g'(a_i) = \prod\limits_{\substack{k=1 \\ k\neq i}}^n (a_i-a_k) \)。
	多项式 \( F(x) \) 的次数为 \( n-1 \)，且有 \( n \) 个不同的根 \( a_1,\dots,a_n \)。一个次数小于 \( n \) 的多项式若有 \( n \) 个不同的根，则必为零多项式，故 \( F(x) = 0 \)，恒等式得证。
	
	再证 \( f \) 关于 \( g \) 的余数，可以写成：
	\[
	\sum_{i=1}^n \frac{f(a_i)g(x)}{(x-a_i)g'(a_i)} \tag{1}
	\]
	
	设 \( f = qg + r \)，其中 \( \deg g > \deg r \)。考虑多项式：
	\[
	F_1(x) = r(x) - \sum_{i=1}^n \frac{f(a_i)g(x)}{(x-a_i)g'(a_i)}
	\]
	则 \( \deg F_1 < n \)，且 \( F_1(a_1)=0, \dots, F_1(a_n)=0 \)。因此 \( F_1(x)=0 \)。
	
	观察式 (1) 的形式，不难取
	\[
	L(x) = \sum_{i=1}^n \frac{b_i g(x)}{(x-a_i)g'(a_i)}
	\]
	它满足 \( L(a_j) = b_j \)，对所有 \( 1 \le j \le n \)。
	
	若 \( H \in \F[x] \) 且 \( \deg H \le n-1 \)，且 \( H(a_j) = b_j \)，则 \( H(a_j) = L(a_j) \)。于是 \( H-L \) 有超过 \( n-1 \) 个根，故 \( H(x) \equiv L(x) \)，即 \( L(x) \) 唯一。
	
\end{proof}
 下面叙述 Newton 公式。令 \( S_k = x_1^k + \dots + x_n^k \)，且 \( S_0 = n \)。
 \begin{lemma}
 	设
 	\[
 	f(x) = (x-x_1)\cdots(x-x_n) = x^n - \sigma_1 x^{n-1} + \dots + (-1)^n \sigma_n
 	\]
 	则
 	\[
 	x^{k+1} f'(x) = (S_0 x^k + S_1 x^{k-1} + \dots + S_k) f(x) + g(x)
 	\]
 	其中 \( \deg g < n \)
 \end{lemma}
\begin{proof}
	 我们知道
	\[
	f'(x) = \sum_{i=1}^n \frac{f(x)}{x-x_i}
	\]
	因此
	\[
	\begin{aligned}
		x^{k+1} f'(x) &= \sum_{i=1}^n \frac{x^{k+1} f(x)}{x-x_i} \\
		&= \sum_{i=1}^n \frac{(x^{k+1} - x_i^{k+1} + x_i^{k+1})}{x-x_i} f(x) \\
		&= f(x) \sum_{i=1}^n \frac{x^{k+1} - x_i^{k+1}}{x-x_i} + g(x)
	\end{aligned}
	\]
	其中 \( g(x) = \sum\limits_{i=1}^n \dfrac{x_i^{k+1} f(x)}{x-x_i} \)，且 \( \deg g < n \)。
	
	而
	\[
	\sum_{i=1}^n \frac{x^{k+1} - x_i^{k+1}}{x-x_i} = \sum_{i=1}^n (x^k + x_i x^{k-1} + \dots + x_i^k) = S_0 x^k + S_1 x^{k-1} + \dots + S_k	\]
	于是
	\[
	x^{k+1} f'(x) = (S_0 x^k + \dots + S_k) f(x) + g(x)
	\]
	
	
\end{proof}

\begin{proposition}（Newton 公式）
	记号同上。
	
	若 \(k \le n-1\)，则
	\[
	S_k - S_{k-1}\sigma_1 + S_{k-2}\sigma_2 - \dots + (-1)^{k-1}S_1\sigma_{k-1} + (-1)^k k \sigma_k = 0
	\]
	
	若 \(k \ge n\)，则
	\[
	S_k - S_{k-1}\sigma_1 + S_{k-2}\sigma_2 - \dots + (-1)^n S_{k-n}\sigma_n = 0
	\]
\end{proposition}
\begin{proof}
	对 \(f(x) = x^n - \sigma_1 x^{n-1} + \sigma_2 x^{n-2} - \dots + (-1)^n\sigma_n\) 求导，乘以 \(x^{k+1}\)，得到：
	\[
	x^{k+1}f'(x) = nx^{n+k} - (n-1)\sigma_1 x^{n+k-1} + (n-2)\sigma_2 x^{n+k-2} + \dots + (-1)^{n-1}\sigma_{n-1}x^{k+1}
	\]
	由前面的引理，我们有：
	\[
	x^{k+1}f'(x) = (S_0 x^k + S_1 x^{k-1} + \dots + S_k)f(x) + g(x)
	\]
	比较两边多项式的系数：
	
	当 \(k \le n-1\) 时，得到：
	\[
	S_k - S_{k-1}\sigma_1 + S_{k-2}\sigma_2 - \dots + (-1)^{k-1}S_1\sigma_{k-1} + (-1)^k k \sigma_k = 0
	\]
	
	当 \(k \ge n\) 时，得到：
	\[
	S_k - S_{k-1}\sigma_1 + \dots + (-1)^n S_{k-n}\sigma_n = 0
	\]
\end{proof}


根据 Newton 公式，可以得到前几个幂和：
\[
\begin{aligned}
	S_1 &= \sigma_1 \\
	S_2 &= \sigma_1^2 - 2\sigma_2 \\
	S_3 &= \sigma_1^3 - 3\sigma_1\sigma_2 + 3\sigma_3 \\
	S_4 &= \sigma_1^4 - 4\sigma_1^2\sigma_2 + 4\sigma_1\sigma_3 + 2\sigma_2^2 - 4\sigma_4
\end{aligned}
\]
\section{有价值的习题}

\begin{problem}
使用欧几里得算法求正整数
	\(n=57,128,567,2001\) 的二进制展开。也就是说，求整数 \(a_i=0,1\) 以及 \(a_d=1\)，使得
	\[
	n=a_0+a_1\cdot 2+a_2\cdot 2^2+\cdots+a_d\cdot 2^d.
	\]
	记 \(n=(a_d\cdots a_2a_1a_0)_2\)。例如，
	\[
	2=(10)_2,\quad 3=(11)_2,\quad 15=(1111)_2,\quad
	49=(10010101)_2,\quad 257=(100000001)_2.
	\]
\end{problem}

\begin{solution}
	\[
	\begin{aligned}
		57&=(111001)_2, & 128&=(10000000)_2,\\
		567&=(1000110111)_2, & 2001&=(11111010001)_2.
	\end{aligned}
	\]
	

\end{solution}
这里涉及把十进制数转化为二进制数。
算法如下
\begin{enumerate}
	\item[(i)] 设 \(p>1\) 是整数，则 \(p\)-进制数与十进制数的关系为
	\[
	(a_{n-1}a_{n-2}\cdots a_1a_0)_p
	=a_{n-1}p^{n-1}+a_{n-2}p^{n-2}+\cdots+a_1p+a_0,
	\]
	其中 \(0\le a_i<p\) 为整数，且 \(a_{n-1}\ne0\)。
	
	\item[(ii)] 给定整数 \(p>1\)，\(a>1\) 是任意整数（默认是十进制数）。重复使用带余除法得
	\[
	\begin{aligned}
		a
		&=a_1p+r_0\\
		&=(a_2p+r_1)p+r_0
		=a_2p^2+r_1p+r_0\\
		&=(a_3p+r_2)p^2+r_1p+r_0
		=a_3p^3+r_2p^2+r_1p+r_0\\
		&=\cdots\\
		&=r_{n-1}p^{n-1}+r_{n-2}p^{n-2}+\cdots+r_1p+r_0,
	\end{aligned}
	\]
	其中 \(0\le r_j<p\)，\(0\le j\le n-1\)，且 \(r_{n-1}\ne0\)。于是
	\[
	a=(r_{n-1}r_{n-2}\cdots r_2r_1r_0)_p
	\]
	就是十进制数 \(a\) 的 \(p\)-进制表达式。
\end{enumerate}
\begin{problem}
设 \(a_1,a_2,\ldots,a_n\) 和 \(b_1,b_2,\ldots,b_n\) 均为整数。
	\begin{enumerate}
		\item 设 \(d\) 为整数，证明
		\[
		(da_1,da_2,\ldots,da_n)=|d|(a_1,a_2,\ldots,a_n).
		\]
		\item 设
		\[
		(a_1,a_2,\ldots,a_n)=(b_1,b_2,\ldots,b_n)=1,
		\]
		并且对某个有理数 \(c\) 及所有 \(1\le i\le n\)，都有 \(b_i=ca_i\)。证明 \(c=\pm1\)。
	\end{enumerate}
\end{problem}

\begin{solution}
	\textbf{注：}原题中应加上 \(a_i\) 和 \(b_i\) 都不全为 \(0\) 的条件。
	
	\begin{enumerate}
		\item 设
		\[
		x=(da_1,da_2,\ldots,da_n),\qquad y=(a_1,a_2,\ldots,a_n).
		\]
		先证 \(|d|y\) 是 \(da_1,\ldots,da_n\) 的公因子。事实上，由 \(y\mid a_i\) 得
		\[
		|d|y\mid da_i,\qquad 1\le i\le n.
		\]
		其次，若 \(t\) 是 \(da_1,\ldots,da_n\) 的任意公因子，由 贝祖 等式，存在 \(u_i\in\Z\) 使
		\[
		\sum_{i=1}^n u_ia_i=y.
		\]
		因而
		\[
		\sum_{i=1}^n u_i(da_i)=dy,
		\]
		故 \(t\mid dy\)，从而 \(t\mid |d|y\)。所以 \(x=|d|y\)。
		
		\item 写 \(c=\dfrac{s}{r}\)，其中 \(r,s\in\Z\) 且 \(r\ne0\)。由 \(b_i=ca_i\) 得
		\[
		rb_i=sa_i,
		\qquad 1\le i\le n.
		\]
		由第 (1) 问及已知条件，
		\[
		(rb_1,\ldots,rb_n)=|r|(b_1,\ldots,b_n)=|r|,
		\]
		\[
		(sa_1,\ldots,sa_n)=|s|(a_1,\ldots,a_n)=|s|.
		\]
		两组整数逐项相等，故 \(|r|=|s|\)，于是
		\[
		c=\frac{s}{r}=\pm1.
		\]
	\end{enumerate}
\end{solution}

\begin{problem}
 设
	\[
	a=\pm p_1^{r_1}p_2^{r_2}\cdots p_n^{r_n},\qquad
	b=\pm p_1^{s_1}p_2^{s_2}\cdots p_n^{s_n}
	\]
	是 \(a\) 与 \(b\) 的素因数分解，其中 \(r_i\) 与 \(s_i\) 均为非负整数。证明
	\[
	(a,b)=p_1^{t_1}p_2^{t_2}\cdots p_n^{t_n},
	\qquad t_i=\min(r_i,s_i),\quad 1\le i\le n.
	\]
\end{problem}

\begin{solution}
	设
	\[
	d=p_1^{t_1}\cdots p_n^{t_n}.
	\]
	由于 \(t_i=\min(r_i,s_i)\)，所以 \(d\mid a\) 且 \(d\mid b\)。另一方面，设 \(m\) 是 \(a,b\) 的任意公因子。由算术基本定理，
	\[
	m=\pm p_1^{t_1'}\cdots p_n^{t_n'}.
	\]
	由 \(m\mid a\) 和 \(m\mid b\) 得
	\[
	t_i'\le r_i,\qquad t_i'\le s_i,
	\]
	故 \(t_i'\le t_i\)，从而 \(m\mid d\)。因此 \(d=(a,b)\)。
\end{solution}

\begin{problem}
给定正整数 \(m\)。若整数 \(a\) 与 \(b\) 满足 \(m\mid a-b\)，则称 \(a\) 与 \(b\) 模 \(m\) 同余，记作
	\(a\equiv b\pmod m\)。例如，
	\[
	2007\equiv7\pmod{2000},\qquad 2007\equiv-1\pmod{2008}.
	\]
	\begin{enumerate}
		\item 证明：\(a\equiv b\pmod m\) 当且仅当 \(a\) 与 \(b\) 除以 \(m\) 所得的最小非负余数相同。
		\item 证明：\(a\equiv0\pmod m\) 当且仅当 \(m\mid a\)。
		\item 若 \(a\equiv a_1\pmod m\) 且 \(b\equiv b_1\pmod m\)，证明
		\[
		a+b\equiv a_1+b_1\pmod m,
		\qquad ab\equiv a_1b_1\pmod m.
		\]
		\item 若 \((a,m)=1\)，证明存在整数 \(b\)，使得
		\[
		ab\equiv1\pmod m.
		\]
		此时称同余方程 \(ax\equiv1\pmod m\) 有解 \(x\equiv b\pmod m\)。
		\item 求整数 \(b\)，使得 \(3b\equiv1\pmod{2006}\)。
	\end{enumerate}
\end{problem}

\begin{solution}
	\begin{enumerate}
		\item “若”部分：设
		\[
		a=a_1m+r,
		\qquad b=b_1m+r
		\]
		是带余除法所得的表示，则
		\[
		a-b=(a_1-b_1)m,
		\]
		故 \(m\mid(a-b)\)。
		
		“仅若”部分：设
		\[
		a=a_1m+r_1,
		\qquad b=b_1m+r_2,
		\qquad 0\le r_1,r_2<m.
		\]
		若 \(a\equiv b\pmod m\)，则 \(m\mid(r_1-r_2)\)。由于 \(|r_1-r_2|<m\)，必有 \(r_1-r_2=0\)。
		
		\item 在第 (1) 问中取 \(b=0\) 即得。
		
		\item 由假设可写成
		\[
		a-a_1=km,
		\qquad b-b_1=lm,
		\qquad k,l\in\Z.
		\]
		于是
		\[
		(a+b)-(a_1+b_1)=(k+l)m,
		\]
		故 \(a+b\equiv a_1+b_1\pmod m\)。另外，
		\[
		ab-a_1b_1=(a-a_1)b+a_1(b-b_1)=(bk+a_1l)m,
		\]
		故 \(ab\equiv a_1b_1\pmod m\)。
		
		\item 由贝祖等式，存在 \(b,v\in\Z\)，使得
		\[
		ba+vm=(a,m)=1.
		\]
		于是 \(m\mid(ba-1)\)，从而 \(ba\equiv1\pmod m\)。
		
		\item 由于 \((3,2006)=1\)，该同余方程有解。事实上，
		\[
		1=669\cdot3-2006,
		\]
		所以
		\[
		669\cdot3\equiv1\pmod{2006}.
		\]
		因此可取 \(b=669\)。
	\end{enumerate}
\end{solution}

\begin{problem}
若 \(a^n-1\) 是素数，\(a>1\) , \(n>1\)证明 \(a=2\) 且 \(n\) 是素数。

\end{problem}
	形如
\[
M_p=2^p-1
\]
\phantomsection\label{p1:concept:c1-l1974}的素数称为\textbf{梅森素数}。
\begin{solution}
	
	
	注意到
	\[
	a^n-1=(a-1)(a^{n-1}+a^{n-2}+\cdots+a+1).
	\]
	若 \(a>2\)，则上式给出 \(a^n-1\) 的非平凡分解，与其为素数矛盾，故 \(a=2\)。
	
	若 \(n=pq\) 不是素数，其中 \(p,q>1\)，则
	\[
	2^n-1=(2^p)^q-1
	=(2^p-1)\bigl((2^p)^{q-1}+\cdots+2^p+1\bigr),
	\]
	仍与 \(2^n-1\) 为素数矛盾。因此 \(n\) 必为素数。
	
	
\end{solution}
\textbf{注：}这只给出了 \(a^n-1\) 为素数的必要条件。形如 \(2^p-1\) 的素数称为 梅森素数。关于 梅森素数的基本问题是：是否有无穷多个 梅森素数。
\begin{problem}
若 \(a^n+1\) 是素数，\(a>1\) , \(n>1\)证明 \(a\) 是偶数且 \(n\) 是 \(2\) 的幂。
	形如
	\[
	F_n=2^{2^n}+1
	\]
	\phantomsection\label{p1:concept:c1-l2000}的素数称为\textbf{费马素数}。例如，当 \(n=0,1,2,3,4\) 时，\(F_n\) 均为素数。欧拉发现
	\[
	F_5=641\cdot6700417
	\]
	该分解于 1732 年得到。目前尚不知道费马素数是否有无穷多个。
\end{problem}

\begin{solution}

	
	若 \(a\) 是奇数，则 \(a^n+1>2\) 是偶数，与其为素数矛盾，故 \(a\) 为偶数。
	
	若 \(n\) 不是 \(2\) 的幂，则由算术基本定理，存在奇素数 \(p\) 使 \(p\mid n\)。写 \(n=pq\)，则
	\[
	\begin{aligned}
		a^n+1
		&=a^{pq}+1=(a^q)^p+1\\
		&=(a^q+1)\bigl((a^q)^{p-1}-(a^q)^{p-2}+\cdots-a^q+1\bigr),
	\end{aligned}
	\]
	这里用了奇数 \(p\) 时的恒等式
	\[
	x^p+1=(x+1)(x^{p-1}-x^{p-2}+\cdots-x+1).
	\]
	这与 \(a^n+1\) 为素数矛盾，故 \(n\) 必为 \(2\) 的幂。
	
	\textbf{注：}这只给出了 \(a^n+1\) 为素数的必要条件。若 \(F_k=2^{2^k}+1\) 是素数，则称 \(F_k\) 为 费马素数。目前只知道 \(F_0,F_1,F_2,F_3,F_4\) 是 费马素数。关于 费马素数的基本问题是：是否有无穷多个 费马素数。
\end{solution}

\begin{problem}
 假设
	\[
	\frac{f(x)}{(f(x),g(x))}
	\quad\text{和}\quad
	\frac{g(x)}{(f(x),g(x))}
	\]
	均不是常数。
	\begin{enumerate}
		\item 证明存在多项式 \(u_0(x),v_0(x)\)，满足
		\[
		\deg u_0(x)<\deg\frac{g(x)}{(f(x),g(x))},
		\qquad
		\deg v_0(x)<\deg\frac{f(x)}{(f(x),g(x))},
		\]
		使得
		\[
		(f(x),g(x))=u_0(x)f(x)+v_0(x)g(x).
		\]
		\item 若
		\[
		(f(x),g(x))=u(x)f(x)+v(x)g(x),
		\]
		证明 \(u(x),v(x)\) 必具有如下形式：
		\[
		u(x)=u_0(x)+\frac{g(x)}{(f(x),g(x))}h(x),
		\qquad
		v(x)=v_0(x)-\frac{f(x)}{(f(x),g(x))}h(x)
		\]
		其中 \(h(x)\) 为某个多项式。
	\end{enumerate}
\end{problem}

\begin{solution}
记 \(d(x)=(f(x),g(x))\)。
	
	\begin{enumerate}
		\item 等价地，要找满足次数条件的 \(u_0,v_0\)，使
		\begin{equation}
			u_0(x)\frac{f(x)}{d(x)}+v_0(x)\frac{g(x)}{d(x)}=1.
			\tag{$*$}
		\end{equation}
		由于
		\[
		\left(\frac{f(x)}{d(x)},\frac{g(x)}{d(x)}\right)=1,
		\]
		由 贝祖 等式，存在 \(u,v\in\F[x]\) 使
		\begin{equation}
			u(x)\frac{f(x)}{d(x)}+v(x)\frac{g(x)}{d(x)}=1.
			\tag{1}
		\end{equation}
		若已经满足所需次数条件，则取 \(u_0=u,v_0=v\)。否则不妨设
		\[
		\deg u(x)\ge\deg\frac{g(x)}{d(x)}.
		\]
		对 \(u(x)\) 作带余除法：
		\begin{equation}
			u(x)=\frac{g(x)}{d(x)}q(x)+r(x),
			\qquad r(x)=0\ \text{或}\ \deg r(x)<\deg\frac{g(x)}{d(x)}.
			\tag{2}
		\end{equation}
		若 \(r(x)=0\)，由 (1)、(2) 得
		\[
		\frac{g(x)}{d(x)}
		\left(\frac{f(x)}{d(x)}q(x)+v(x)\right)=1,
		\]
		这与 \(g(x)/d(x)\) 非常数矛盾。故
		\begin{equation}
			\deg r(x)<\deg\frac{g(x)}{d(x)}.
			\tag{3}
		\end{equation}
		将 (2) 代入 (1)，得
		\begin{equation}
			r(x)\frac{f(x)}{d(x)}+\bigl(v(x)+q(x)\bigr)\frac{g(x)}{d(x)}=1.
			\tag{4}
		\end{equation}
		取
		\[
		u_0(x)=r(x),
		\qquad v_0(x)=v(x)+q(x).
		\]
		由 (3) 得到 \(u_0\) 的次数条件；比较 (4) 两边的次数，得到
		\[
		\deg v_0(x)<\deg\frac{f(x)}{d(x)}.
		\]
		故 \(u_0,v_0\) 满足要求。
		
		\item 由已知等式除以 \(d(x)\)，得
		\begin{equation}
			u(x)\frac{f(x)}{d(x)}+v(x)\frac{g(x)}{d(x)}=1.
			\tag{5}
		\end{equation}
		而第 (1) 问构造的 \(u_0,v_0\) 满足
		\begin{equation}
			u_0(x)\frac{f(x)}{d(x)}+v_0(x)\frac{g(x)}{d(x)}=1.
			\tag{6}
		\end{equation}
		两式相减，得
		\begin{equation}
			\bigl(u(x)-u_0(x)\bigr)\frac{f(x)}{d(x)}
			+\bigl(v(x)-v_0(x)\bigr)\frac{g(x)}{d(x)}=0.
			\tag{7}
		\end{equation}
		于是
		\[
		\frac{g(x)}{d(x)}
		\mid
		\left(\bigl(u(x)-u_0(x)\bigr)\frac{f(x)}{d(x)}\right).
		\]
		又因 \(f(x)/d(x)\) 与 \(g(x)/d(x)\) 互素，故
		\[
		\frac{g(x)}{d(x)}\mid u(x)-u_0(x).
		\]
		因而存在 \(h_1(x)\in\F[x]\) 使
		\[
		u(x)=u_0(x)+\frac{g(x)}{d(x)}h_1(x).
		\]
		同理，存在 \(h_2(x)\in\F[x]\) 使
		\[
		v(x)=v_0(x)+\frac{f(x)}{d(x)}h_2(x).
		\]
		代入 (7) 得 \(h_1(x)+h_2(x)=0\)。取
		\[
		h(x)=h_1(x)=-h_2(x),
		\]
		即得
		\[
		u(x)=u_0(x)+\frac{g(x)}{d(x)}h(x),
		\qquad
		v(x)=v_0(x)-\frac{f(x)}{d(x)}h(x).
		\]
	\end{enumerate}
\end{solution}

\begin{problem}
\phantomsection\label{p1:concept:c1-l2164}若 \(f(x)\mid m(x)\) 且 \(g(x)\mid m(x)\)，则称 \(m(x)\) 是 \(f(x)\) 与 \(g(x)\) 的公倍式。定义 \(f(x)\) 与 \(g(x)\) 的最小公倍式 \([f(x),g(x)]\) 为次数最小的首一公倍式。对有限个多项式的最小公倍式递归定义为
	\[
	[f_1(x),\ldots,f_{n-1}(x),f_n(x)]
	=[[f_1(x),\ldots,f_{n-1}(x)],f_n(x)],
	\qquad n>2.
	\]
	\begin{enumerate}
		\item 设 \(f(x)\) 与 \(g(x)\) 均为首一多项式，证明
		\[
		[f(x),g(x)]=\frac{f(x)g(x)}{(f(x),g(x))}.
		\]
		\item 若 \(m(x)\) 是 \(f_1(x),\ldots,f_n(x)\) 的公倍式，证明
		\[
		[f_1(x),\ldots,f_n(x)]\mid m(x).
		\]
	\end{enumerate}
\end{problem}

\begin{solution}
	\begin{enumerate}
		\item 令
		\[
		d(x)=(f(x),g(x)),
		\qquad M(x)=\frac{f(x)g(x)}{d(x)}.
		\]
		由于 \(d\mid f\) 且 \(d\mid g\)，有
		\[
		M(x)=g(x)\frac{f(x)}{d(x)}=f(x)\frac{g(x)}{d(x)},
		\]
		故 \(M\) 是 \(f,g\) 的公倍式。
		
		设 \(h(x)\) 是任意公倍式，可写成
		\[
		h(x)=f(x)h_1(x)=g(x)h_2(x).
		\]
		于是
		\[
		\frac{f(x)}{d(x)}h_1(x)=\frac{g(x)}{d(x)}h_2(x).
		\]
		由于
		\[
		\left(\frac{f(x)}{d(x)},\frac{g(x)}{d(x)}\right)=1,
		\]
		可得 \(f(x)/d(x)\mid h_2(x)\)。因此
		\[
		M(x)=g(x)\frac{f(x)}{d(x)}\mid g(x)h_2(x)=h(x).
		\]
		因此 \(M\) 是次数最小的首一公倍式，即 \(M=[f,g]\)。
		
		\item 由第 (1) 问知 \([f_1,f_2]\mid m\)。又因 \(f_3\mid m\)，将第 (1) 问用于 \([f_1,f_2]\) 与 \(f_3\)，得
		\[
		[f_1,f_2,f_3]=[[f_1,f_2],f_3]\mid m.
		\]
		重复上述过程即得
		\[
		[f_1,\ldots,f_n]\mid m.
		\]
	\end{enumerate}
\end{solution}

由此得到如下推论。

\begin{corollary}
	设
	\[
	f(x)=p_1(x)^{r_1}\cdots p_k(x)^{r_k},
	\qquad
	g(x)=p_1(x)^{s_1}\cdots p_k(x)^{s_k},
	\]
	其中 \(p_1(x),\ldots,p_k(x)\) 两两不同且不可约，则
	\[
	[f(x),g(x)]
	=p_1(x)^{\max(r_1,s_1)}\cdots p_k(x)^{\max(r_k,s_k)}.
	\]
\end{corollary}

\begin{proof}
	由
	\[
	(f(x),g(x))
	=p_1(x)^{\min(r_1,s_1)}\cdots p_k(x)^{\min(r_k,s_k)}
	\]
	及
	\[
	[f(x),g(x)]=\frac{f(x)g(x)}{(f(x),g(x))}
	\]
	立即得到结论。
\end{proof}

\begin{corollary}
	设 \(f(x),g(x)\in\F[x]\) 都是非零多项式。多项式 \(m(x)\in\F[x]\) 是 \(f(x),g(x)\) 的最小公倍式，当且仅当
	\begin{enumerate}
		\item \(m(x)\) 是 \(f(x),g(x)\) 的公倍式；
		\item 对任意公倍式 \(h(x)\)，都有 \(m(x)\mid h(x)\)。
	\end{enumerate}
	特别地，\([f(x),g(x)]\mid h(x)\)。
\end{corollary}

\begin{proof}
	``\(\Leftarrow\)'' 由最小公倍式的定义立即得到。
	
	``\(\Rightarrow\)'' 设 \([f,g]\) 是按次数定义的最小公倍式。对任意公倍式 \(h\)，作带余除法
	\[
	h(x)=[f(x),g(x)]q(x)+r(x),
	\qquad r(x)=0\ \text{或}\ \deg r<\deg[f,g].
	\]
	由 \(f\mid h\)、\(g\mid h\)、\(f\mid[f,g]\)、\(g\mid[f,g]\)，得 \(f\mid r\)、\(g\mid r\)。若 \(r\ne0\)，则 \(r\) 是次数更小的公倍式，矛盾。因此 \(r=0\)，即 \([f,g]\mid h\)。
\end{proof}

\begin{problem}
（在有理数域上分解因式的 Kronecker算 法。）说明如何找出整系数多项式 \(f(x)\) 在有理数域上的全部因式。
\end{problem}

\begin{solution}
	\textbf{算法原理：}不妨设整系数多项式 \(f(x)\) 是本原多项式。假设 \(f(x)\) 在有理数域上有一个次数 \(\le s\) 的因式。由 高斯 引理，\(f(x)\) 必有一个次数 \(\le s\) 的整系数因式 \(g(x)\)，且 \(f(x)/g(x)\) 也是整系数多项式。
	
	于是对任意整数 \(a\)，都有
	\[
	g(a)\mid f(a).
	\]
	任取互不相同的 \(s+1\) 个整数 \(a_0,a_1,\ldots,a_s\)，使 \(f(a_i)\ne0\)。那么 \(f(x)\) 有次数 \(\le s\) 的因式的必要条件是：存在次数 \(\le s\) 的整系数多项式 \(g(x)\)，使
	\[
	g(a_i)\mid f(a_i),
	\qquad 0\le i\le s.
	\]
	
	对每个 \(i\)，任取 \(f(a_i)\) 的一个因数 \(b_i\)。由 拉格朗日 插值公式，存在唯一的次数 \(\le s\) 的有理系数多项式 \(h(x)\)，使
	\[
	h(a_i)=b_i,
	\qquad 0\le i\le s.
	\]
	由上述必要条件，\(h(x)\) 可能是 \(f(x)\) 的一个次数 \(\le s\) 的因式。于是，对所有可能的
	\[
	(b_0,b_1,\ldots,b_s)
	\]
	所确定的 \(h(x)\) 逐个检验是否整除 \(f(x)\)，即可得到 \(f(x)\) 的全部次数 \(\le s\) 的整系数因式。若所得候选均不整除 \(f(x)\)，则 \(f(x)\) 没有次数 \(\le s\) 的因式。
	
	改变 \(s\)，即可得到 \(f(x)\) 在有理数域上的所有因式。事实上，若对不超过 \(\deg f/2\) 的最大整数 \(s\)，上述方法只得到 \(0\) 次多项式，则 \(f(x)\) 在 \(\Q\) 上不可约。Kronecker 法实际上是一个“筛选法”。
\end{solution}

\begin{problem}
 \phantomsection\label{p1:concept:c1-l2305}设 \(f(x)\) 是复数域 \(\C\) 上的多项式。将 \(f(x)\) 的每个系数取共轭所得的多项式记为 \(\overline{f}(x)\)。
	\begin{enumerate}
		\item 若 \(f(x)\mid g(x)\)，证明 \(\overline{f}(x)\mid\overline{g}(x)\)。
		\item 若 \(\alpha\) 同时是 \(f(x)\) 与 \(\overline{f}(x)\) 的根，\(\alpha\) 是否必为实数？
	\end{enumerate}
\end{problem}

\begin{solution}
	\begin{enumerate}
		\item 由 \(f(x)\mid g(x)\)，存在 \(h(x)\in\C[x]\)，使得
		\[
		g(x)=f(x)h(x).
		\]
		对两边各系数取共轭，得
		\[
		\overline{g}(x)=\overline{f}(x)\,\overline{h}(x),
		\]
		故 \(\overline{f}(x)\mid\overline{g}(x)\)。
		
		\item 不一定。例如，
		\[
		f(x)=x^2+1=\overline{f}(x)
		\]
		具有非实公共根 \(\alpha=i\)。
	\end{enumerate}
\end{solution}

\begin{problem}
	 设 \(f(x)\) 是实数域 \(\R\) 上的多项式，并且对任意 \(a\in\R\) 都有 \(f(a)\ge0\)。证明存在多项式 \(u(x),v(x)\in\R[x]\)，使得
	\[
	f(x)=u(x)^2+v(x)^2.
	\]
\end{problem}

\begin{solution}
	先考虑 \(\R\) 上任意首一不可约二次多项式
	\[
	x^2+px+q.
	\]
	其判别式满足 \(p^2-4q<0\)，配方得
	\begin{equation}
		x^2+px+q
		=\left(x+\frac p2\right)^2
		+\left(\frac{\sqrt{4q-p^2}}2\right)^2,
		\tag{1}
	\end{equation}
	所以它是两个平方之和。
	
	再注意恒等式
	\[
	(a^2+b^2)(c^2+d^2)=(ac+bd)^2+(ad-bc)^2.
	\]
	因此，平方和的乘积仍是平方和。
	
	在 \(\R[x]\) 中分解 \(f\)。由于 \(f(a)\ge0\)，每个实根的重数都是偶数，故可写成
	\[
	f(x)=a\,g(x)h(x),
	\]
	其中 \(g(x)=g_1(x)^2\) 是实线性因子的偶次幂之积，\(h(x)\) 是首一不可约二次多项式之积。由 (1) 和上面的恒等式，存在 \(u_1,v_1\in\R[x]\) 使
	\[
	h(x)=u_1(x)^2+v_1(x)^2.
	\]
	又因 \(f\ge0\) 且 \(f\not\equiv0\) 时首项系数 \(a>0\)，所以
	\[
	\begin{aligned}
		f(x)
		&=a\,g_1(x)^2\bigl(u_1(x)^2+v_1(x)^2\bigr)\\
		&=\bigl(\sqrt a\,g_1(x)u_1(x)\bigr)^2
		+\bigl(\sqrt a\,g_1(x)v_1(x)\bigr)^2.
	\end{aligned}
	\]
	零多项式的情形显然成立。
\end{solution}

\begin{problem}
	 设 \(f(x)\in\Q[x]\)，\(a,b\in\Q\) 且 \(a\ne0\)。证明：\(f(x)\) 在 \(\Q\) 上不可约，当且仅当 \(f(ax+b)\) 在 \(\Q\) 上不可约。
\end{problem}

\begin{solution}
	由于 \(a\ne0\)，有
	\[
	\deg f(x)=\deg f(ax+b).
	\]
	
	``\(\Rightarrow\)'' 假设 \(f(x)\) 在 \(\Q\) 上不可约。若
	\[
	f(ax+b)=g_1(x)g_2(x),
	\qquad \deg g_i<\deg f,
	\]
	则令
	\[
	h_i(x)=g_i\left(\frac{x-b}{a}\right),
	\qquad i=1,2,
	\]
	得到
	\[
	f(x)=h_1(x)h_2(x),
	\qquad \deg h_i=\deg g_i<\deg f,
	\]
	与 \(f\) 不可约矛盾。因此 \(f(ax+b)\) 不可约。
	
	``\(\Leftarrow\)'' 假设 \(g(x)=f(ax+b)\) 在 \(\Q\) 上不可约。将已经证明的方向应用于 \(g\) 以及线性代换 \(x\mapsto(x-b)/a\)，得到
	\[
	g\left(\frac{x-b}{a}\right)=f(x)
	\]
	在 \(\Q\) 上不可约。
\end{solution}

\begin{problem}
	 设
	\[
	f(x)=a_nx^n+a_{n-1}x^{n-1}+\cdots+a_1x+a_0\in\Z[x].
	\]
	设 \(p\) 为素数，并且
	\[
	p\nmid a_n,
	\qquad p\mid a_i\quad(0\le i\le k-1),
	\qquad p^2\nmid a_0,
	\]
	其中 \(k\) 为正整数。证明 \(f(x)\) 在 \(\Q\) 上有一个次数至少为 \(k\) 的不可约因式。
\end{problem}

\begin{solution}
	本题是 艾森斯坦 判别法的推广。当 \(k=n\) 时，结论给出 \(f(x)\) 有一个次数至少为 \(n\) 的不可约因式，也就是通常的 艾森斯坦 判别法。
	
	若 \(f(x)\) 在 \(\Q\) 上不可约，结论显然成立。以下设 \(f(x)\) 可约。由 高斯 引理，可将其写成整系数不可约多项式之积：
	\[
	f(x)=p_1(x)^{r_1}p_2(x)^{r_2}\cdots p_s(x)^{r_s},
	\]
	其中各 \(p_i(x)\in\Z[x]\) 在 \(\Q\) 上不可约且两两不相同。
	
	反设所有 \(p_i\) 的次数都小于 \(k\)。设 \(p_i(x)\) 的常数项为 \(b_{i0}\)，则
	\[
	a_0=\prod_{i=1}^s b_{i0}^{r_i}.
	\]
	由 \(p\mid a_0\) 且 \(p^2\nmid a_0\)，恰有一个因式的常数项被 \(p\) 整除，而且该因式只出现一次。不妨记为
	\[
	p_1(x)=b_mx^m+\cdots+b_1x+b_0,
	\qquad m<k,
	\]
	其中 \(p\mid b_0\)。令
	\[
	g(x)=p_2(x)^{r_2}\cdots p_s(x)^{r_s}
	=c_lx^l+\cdots+c_1x+c_0.
	\]
	其余因式的常数项都不被 \(p\) 整除，故
	\begin{equation}
		p\nmid c_0.
		\tag{1}
	\end{equation}
	又由 \(p\nmid a_n\) 可知 \(p\nmid b_m\)。结合 \(p\mid b_0\)，存在 \(1\le j\le m<k\)，使
	\begin{equation}
		p\mid b_0,b_1,\ldots,b_{j-1},
		\qquad p\nmid b_j.
		\tag{2}
	\end{equation}
	由 (1)、(2) 及 \(p\) 为素数，
	\begin{equation}
		p\nmid b_jc_0.
		\tag{3}
	\end{equation}
	而 \(x^j\) 的系数为
	\begin{equation}
		a_j=b_0c_j+b_1c_{j-1}+\cdots+b_{j-1}c_1+b_jc_0.
		\tag{4}
	\end{equation}
	由 (2) 可知 (4) 中除最后一项外均被 \(p\) 整除，而由 (3) 最后一项不被 \(p\) 整除。因此 \(p\nmid a_j\)。但 \(j\le m<k\)，这与假设 \(p\mid a_0,a_1,\ldots,a_{k-1}\) 矛盾。
	
	故至少有一个不可约因式的次数不小于 \(k\)。
\end{solution}

\begin{problem}
	 设
	\[
	f(x)=x^n+a_1x^{n-1}+\cdots+a_{n-1}x+a_n
	=(x-x_1)(x-x_2)\cdots(x-x_n).
	\]
	\phantomsection\label{p1:concept:c1-l2482}定义 \(f\) 的判别式为
	\[
	D(f)=\prod_{1\le i<j\le n}(x_i-x_j)^2.
	\]
	容易看出，\(f(x)\) 有重根当且仅当 \(D(f)=0\)。
	\begin{enumerate}
		\item 当 \(n=2\) 时，证明 \(D(f)=a_1^2-4a_2\)。
		\item 当 \(n=3\) 时，证明
		\[
		D(f)=a_1^2a_2^2-4a_2^3-4a_1^3a_3-27a_3^2+18a_1a_2a_3.
		\]
		\item 设 \(n=3\) 且 \(f(x)\in\R[x]\)。证明：若 \(D(f)>0\)，则 \(f(x)\) 有三个互不相同的实根；若 \(D(f)<0\)，则 \(f(x)\) 只有一个实根。
	\end{enumerate}
\end{problem}

\begin{solution}
	\begin{enumerate}
		\item 当 \(n=2\) 时，
		\[
		D(f)=(x_1-x_2)^2=(x_1+x_2)^2-4x_1x_2.
		\]
		由 韦达 定理，\(x_1+x_2=-a_1\)，\(x_1x_2=a_2\)，故
		\[
		D(f)=a_1^2-4a_2.
		\]
		这正是首一二次方程 \(x^2+a_1x+a_2=0\) 的判别式。
		
		\item 设
		\[
		f(x)=x^3+a_1x^2+a_2x+a_3,
		\]
		并记基本对称多项式为
		\[
		\sigma_1=x_1+x_2+x_3,
		\qquad \sigma_2=x_1x_2+x_1x_3+x_2x_3,
		\qquad \sigma_3=x_1x_2x_3.
		\]
		判别式
		\[
		D(f)=(x_1-x_2)^2(x_1-x_3)^2(x_2-x_3)^2
		\]
		是关于 \(x_1,x_2,x_3\) 的对称多项式，其首项指数组为 \((4,2,0)\)。满足条件的指数组为
		\[
		(4,2,0),\quad(4,1,1),\quad(3,3,0),\quad(3,2,1),\quad(2,2,2).
		\]
		由对称多项式基本定理，可设
		\begin{equation}
			D(f)=r_1\sigma_1^2\sigma_2^2+r_2\sigma_1^3\sigma_3
			+r_3\sigma_2^3+r_4\sigma_1\sigma_2\sigma_3+r_5\sigma_3^2.
			\tag{1}
		\end{equation}
		取若干组特殊值确定系数：
		
		取 \((x_1,x_2,x_3)=(1,-1,0)\)，得 \(D=4\)，\(\sigma_1=0\)，\(\sigma_2=-1\)，\(\sigma_3=0\)，故
		\begin{equation}
			-r_3=4,
			\qquad r_3=-4.
			\tag{2}
		\end{equation}
		
		取 \((x_1,x_2,x_3)=(1,1,-2)\)，得 \(D=0\)，\(\sigma_1=0\)，\(\sigma_2=-3\)，\(\sigma_3=-2\)，故
		\begin{equation}
			-27r_3+4r_5=0,
			\qquad r_5=-27.
			\tag{3}
		\end{equation}
		
		取 \((x_1,x_2,x_3)=(0,2,-1)\)，得 \(D=36\)，\(\sigma_1=1\)，\(\sigma_2=-2\)，\(\sigma_3=0\)，故
		\begin{equation}
			4r_1-8r_3=36,
			\qquad r_1=1.
			\tag{4}
		\end{equation}
		
		取 \((x_1,x_2,x_3)=(-1,1,1)\)，得 \(D=0\)，\(\sigma_1=1\)，\(\sigma_2=-1\)，\(\sigma_3=-1\)。代入 (1)，并利用 (2)--(4)，得
		\begin{equation}
			-r_2+r_4=22.
			\tag{5}
		\end{equation}
		
		取 \((x_1,x_2,x_3)=(1,2,-1)\)，得 \(D=36\)，\(\sigma_1=2\)，\(\sigma_2=-1\)，\(\sigma_3=-2\)。代入 (1)，得
		\[
		4r_1-16r_2-r_3+4r_4+4r_5=36,
		\]
		结合 (2)--(4)，有
		\begin{equation}
			-4r_2+r_4=34.
			\tag{6}
		\end{equation}
		由 (5)、(6) 解得
		\begin{equation}
			r_2=-4,
			\qquad r_4=18.
			\tag{7}
		\end{equation}
		因此
		\begin{equation}
			D(f)=\sigma_1^2\sigma_2^2-4\sigma_1^3\sigma_3-4\sigma_2^3
			+18\sigma_1\sigma_2\sigma_3-27\sigma_3^2.
			\tag{8}
		\end{equation}
		由 韦达 定理，
		\[
		\sigma_1=-a_1,
		\qquad \sigma_2=a_2,
		\qquad \sigma_3=-a_3.
		\]
		代入 (8) 得
		\[
		D(f)=a_1^2a_2^2-4a_1^3a_3-4a_2^3+18a_1a_2a_3-27a_3^2.
		\]
		
		\item 设 \(D(f)>0\)。若 \(f\) 有重根，则 \(D(f)=0\)，矛盾，故各根互不相同。若有非实根，则实系数多项式的虚根成共轭对出现，可写成
		\[
		x_1=c+di,
		\qquad x_2=c-di,
		\qquad x_3\in\R,
		\qquad d\ne0.
		\]
		于是
		\[
		(x_1-x_2)^2=-4d^2<0,
		\]
		而
		\[
		(x_1-x_3)(x_2-x_3)=(c-x_3)^2+d^2>0.
		\]
		故
		\[
		D(f)=(x_1-x_2)^2\bigl((x_1-x_3)(x_2-x_3)\bigr)^2<0,
		\]
		矛盾。因此 \(f\) 有三个互不相同的实根。
		
		再设 \(D(f)<0\)。三次实系数多项式至少有一个实根。若实根多于一个，则由于非实根成共轭对出现，\(f\) 的三个根都为实根。若有重根，则 \(D(f)=0\)；若三个实根互不相同，则
		\[
		D(f)=\prod_{i<j}(x_i-x_j)^2>0.
		\]
		两者均与 \(D(f)<0\) 矛盾。因此 \(f\) 恰有一个实根。
	\end{enumerate}
	
	\textbf{注：}这里给出了三次实系数多项式判别式的应用。
\end{solution}











































































